\documentclass[preprint, 11pt,nonatbib]{elsarticle}
\usepackage[utf8]{inputenc}
\usepackage{todonotes}
\usepackage{amssymb}
\usepackage{amsmath}
\usepackage{amsthm}
\usepackage{tikz-cd}
\usepackage{verbatim}

\usepackage[title]{appendix}


\def\V{{\mathcal V}}

\setlength{\oddsidemargin}{.5in}
\setlength{\evensidemargin}{.5in}
\setlength{\textwidth}{5.5in}

\makeatletter
\let\c@author\relax
\makeatother


\usepackage[
backend=biber,
%style=alphabetic,
sorting=nty,doi=false,url=false
]{biblatex}
\addbibresource{citations.bib}


\DeclareMathOperator{\Dimension}{dim}
\DeclareMathOperator{\Dom}{Val}
\DeclareMathOperator{\Proj}{Proj}
\DeclareMathOperator{\ord}{ord}
\DeclareMathOperator{\Spec}{Spec}
\DeclareMathOperator{\Zar}{Zar}
\DeclareMathOperator{\GreatestCommonDivisor}{gcd}
\DeclareMathOperator{\TranscendenceDegree}{tr. deg.}

\newtheorem{theorem}{Theorem}
\newtheorem{proposition}[theorem]{Proposition}
\newtheorem{corollary}[theorem]{Corollary}
\newtheorem{lemma}[theorem]{Lemma}

\theoremstyle{definition}
\newtheorem{definition}[theorem]{Definition}
\newtheorem{notation}[theorem]{Notation}
\newtheorem{remark}[theorem]{Remark}
\newtheorem{setting}[theorem]{Setting}
\newtheorem{example}[theorem]{Example}
\newtheorem{question}[theorem]{Question}

\numberwithin{theorem}{section}

\begin{document}

\begin{frontmatter}

%% Title, authors and addresses

%% use the tnoteref command within \title for footnotes;
%% use the tnotetext command for theassociated footnote;
%% use the fnref command within \author or \address for footnotes;
%% use the fntext command for theassociated footnote;
%% use the corref command within \author for corresponding author footnotes;
%% use the cortext command for theassociated footnote;
%% use the ead command for the email address,
%% and the form \ead[url] for the home page:
%% \title{Title\tnoteref{label1}}
%% \tnotetext[label1]{}
%% \author{Name\corref{cor1}\fnref{label2}}
%% \ead{email address}
%% \ead[url]{home page}
%% \fntext[label2]{}
%% \cortext[cor1]{}
%% \address{Address\fnref{label3}}
%% \fntext[label3]{}




\title{A connectedness theorem for spaces of valuation rings} 

%\title{Integrally closed local rings and connected sets of valuation rings}



\author[Heinzer]{William Heinzer}
\address[Heinzer]{Department of Mathematics, Purdue University, West
Lafayette, Indiana 47907-1395}
\ead{heinzer@purdue.edu}

\author[Loper]{K.~Alan Loper}
\address[Loper]{Department of Mathematics, Ohio State University -- Newark,  Newark, OH 43055}
\ead{lopera@math.ohio-state.edu}

\author[Olberding]{Bruce Olberding} 
\address[Olberding]{Department of Mathematical Sciences, New Mexico State University,
 Las Cruces, NM 88003-8001}
 \ead{bruce@nmsu.edu}

\author[Toeniskoetter]{Matthew Toeniskoetter} 
\address[Toeniskoetter]{Department of Mathematics and Statistics, Oakland University, Rochester, MI 48309-4479}
 \ead{toeniskoetter@oakland.edu}
 

%% use optional labels to link authors explicitly to addresses:
%% \author[label1,label2]{}
%% \address[label1]{}
%% \address[label2]{}

%\author{}

%\address{}

\begin{abstract}
%% Text of abstract
Let $F$ be a field,  let $D$ be a  
local subring of $F$, and let $\Dom_F(D)$ be the space of valuation rings of $F$ that dominate $D$. We lift Zariski's connectedness theorem for fibers of a projective morphism  to the Zariski-Riemann space of  valuation rings of $F$ by proving that 
 a subring $R$ of $F$ dominating $D$  is local, residually algebraic over $D$ and integrally closed in $F$ if and only if there is a closed and connected subspace $Z$  of $\Dom_F(D)$ such that $R$ is the  intersection of the rings in $Z$.
 Consequently, the intersection of the rings in any closed and connected subset of $\Dom_F(D)$ is a local ring. 
 In proving this, we also prove a converse to 
 Zariski's connectedness theorem. Our results do not  require   the rings involved to be Noetherian. 
%
%\subseteq R$ be   
% subrings of $F$ with $D$ local.  We give a topological criterion for when a subspace of the space of vala
%We prove that
%a subring $R$ of $F$ dominating $D$  is local, residually algebraic over $D$ and integrally closed in $F$ if and only if there is a closed and connected subspace $Z$  of valuation rings of $F$ that dominate $D$  such that $R$ is the  intersection of the rings in $Z$.


\end{abstract}



\begin{keyword}
Connectedness theorem \sep
local ring \sep
valuation ring \sep
Zariski-Riemann space 
%regular local ring \sep  quadratic transform \sep   quadratic tree \sep 
  

\MSC 13A18 \sep 13H99

%13A15 \sep 13C05 \sep 13E05 \sep 13H15

%% keywords here, in the form: keyword \sep keyword

%% PACS codes here, in the form: \PACS code \sep code

%% MSC codes here, in the form: \MSC code \sep code
%% or \MSC[2008] code \sep code (2000 is the default)

\end{keyword}

\end{frontmatter}




\section{Introduction}

In one of its traditional forms, Zariski's Connectedness Theorem asserts that 
if 
$D$ is a Noetherian local
 subring of a field $F$ that is integrally closed in $F$  
and $f:X \rightarrow \Spec(D)$ is a dominant projective morphism of  integral schemes such that $X$ has
 function field contained in  $F$, then  the fiber of $f$ over the closed point of $\Spec(D)$ is closed and connected (cf.~\cite[Corollary~11.3, p.~279]{MR0463157}). 
 That $D$ is local and integrally closed is needed here, and part of the purpose of this article is to examine the relationship between these two properties and the connectedness of the closed fibers of projective morphisms. We do this by fixing the field $F$ and its subring $D$ and varying the projective morphisms $X \rightarrow \Spec(D)$.   
Taking the inverse limit of all such projective morphisms   for which $X$ has function field contained in $F$ results in the Zariski-Riemann space of valuation rings of $F$ that contain $D$. It is in this space that we consider the relationship between local rings and connectedness. 
 


% Thus  if $D$ is a normal Noetherian local  domain, and $f:X \rightarrow \Spec D$ is a blowup of an ideal of $D$, the fiber  $f^{-1}({\mathfrak{m}}_D)$ over the maximal ideal ${\mathfrak{m}}_D$ of $D$ is connected. That $D$ is both local and normal is necessary here for the validity of the theorem, and part of the goal of this article is to understand better this necessity. One way we do this is by fixing a field $F$ containing the local domain $D$ and varying the projective morphisms. Taking the inverse limit of  all projective morphisms $X \rightarrow \Spec D$ such that $X$ has function field contained in $F$ results in the Zariski-Riemann space of valuation rings of $F$ that contain $D$.  It is in this space that we examine the connection between local rings and connectness.  


In more detail,
a {\it valuation ring}  of a field $F$ is a subring $V$ of $F$ such that for all $0 \ne x \in F$, $x \in V$ or $x^{-1} \in V$. 
If $D$ is a subring of $F$, then 
 the {\it Zariski-Riemann space} $\Zar_F(D)$ of $F$ is the space of valuation rings of $F$ that contain $D$ as a subring, where $\Zar_F(D)$ is endowed with the Zariski topology (see Section 2). If also $D$ is local, then $\Dom_F(D)$ is the subspace  of $\Zar_F(D)$ consisting of the valuation rings $V$    that {\it dominate} $D$, i.e.\ the maximal ideal of $V$ contains the maximal ideal of $D$. We show that Zariski's Connectedness Theorem lifts to $\Dom_F(D)$ when $D$ is a local ring that is integrally closed in $F$. Thanks to a recent non-Noetherian version of Stein factorization in The Stacks Project \cite[{Tag~03H2}]{stacks-project}, we do not need to assume $D$ is Noetherian; see Theorem~\ref{first direction}.  
 We also prove a converse to this and give a purely topological criterion for when an intersection of valuation rings is a local ring that is  residually algebraic over $D$. The main purpose of this note is to prove:

\begin{theorem}[Valuative Connectedness Theorem]\label{ValuativeConnectednessTheorem}
Let $F$ be a field, and let $D$ be a local subring of $F$.
A subring $R$ of $F$ dominating $D$ is local, residually algebraic over $D$, and integrally closed in $F$ if and only if there is a closed and connected subspace $Z$ of $\Dom_F(D)$ such that $R$ is the intersection of the rings in $Z$.
\end{theorem} 

 
%Every integrally closed domain is the intersection of the  valuation rings of its quotient field that contain it. In this article we examine the subspaces of $\Zar(F)$ whose valuation rings intersect to a local ring, and we show  
%that  connectedness is a decisive topological feature for determining whether an intersection of valuation rings is local. 

%The issue of when valuation rings intersect to a local ring is closely related to that of the connectedness of fibers of projective  morphisms of schemes, such as arises in . 





Using the fact that the continuous image of a connected set is connected, along with the valuative criterion for properness, one can show that Theorem~\ref{ValuativeConnectednessTheorem} implies Zariski's Connectedness Theorem.
Thus one cannot expect to bypass the difficulty of proving Zariski's theorem by giving a direct proof of Theorem~\ref{ValuativeConnectednessTheorem}.
 
 
 %can be deduced from  \cite[Corollary 11.5, p.~280]{MR0463157}.)
%
%$f:X \rightarrow Y$ is a birational projective morphism of Noetherian integral schemes with $Y$ normal, then for each point $y \in Y$, the set $f^{-1}(y)$ is connected; see {\bf Hartshorne, Corollary 11.4, p.~280.} The theorem can be rephrased in local form to assert that
%For example, if  $D$ is a  normal local domain with maximal ideal ${\mathfrak{m}}$ and
%$X= {\rm Proj} \: D[It]$ is the blow-up along a finitely generated ideal $I$ of $D$, then the 
%fiber $f^{-1}({\mathfrak{m}})$ over ${\mathfrak{m}}$ of the map $f:X \rightarrow \Spec D$  is    connected.
%The stalks of the scheme Proj$(D[It])$ can be viewed as local rings that contain $D$ and have the same quotient field as $D$. 
%We will show  this connectedness theorem lifts to the space of valuation rings  of $F$ in the following form (see Theorem~\ref{Noetherian local}): 


 
%As a corollary, we obtain that Zariski's connectedness theorem
%holds also for non-Noetherian choices of $D$, if $D$ is a directed limit of Noetherian local subrings integrally closed in $F$. (We do not know whether the restriction to choices of  $D$ that are such direct limits is necessary; see Section 4.)
 %In particular, if $R$ is a local overring of an analytically unramified Noetherian local domain, then both the valuative connectedness theorem and Zariski's connectedness theorem hold for $R$ viewed as a subring of its quotient field; see Corollaries~\ref{Nagata ring} and~\ref{au}. 
 
 

%More generally,  we examine the relationship between integrally closed local domains $D$ and closed and connected sets of valuation rings that intersect to  $D$.  
%%In Theorem~\ref{is local} we prove a kind of converse of this theorem by showing that if $D$ is a local subring of a field $F$ and 
%%$Z$ is a closed and connected subset of the space of valuation rings of $F$   that dominate $D$, then the intersection of valuation rings in $Z$ is local.  The assumption that $Z$ is closed is important here, as an example in Section~4 shows. 
%We also prove in Corollary~\ref{au} a version of the valuative connectedness theorem that  works for not-necessarily-Noetherian  overrings of Noetherian local domains: If $R$ is an integrally closed local overring of an analytically unramified  Noetherian local domain, then the set of valuation rings of $F$  that dominate $R$ is connected.

%We will work with spaces of local subrings of a field under two topologies, the Zariski and patch topology, both of which are defined in the next section. These topologies are spectral, in the sense defined in the appendix. To streamline some of the arguments as well as pinpoint    topological arguments that underlie our approach, we collect 
%several needed observations about spectral  spaces 
%in an appendix of the paper. 


%A stronger version of this theorem is proved 
%in Theorem~\ref{two dim} for the case of an overring of a two-dimensional regular local ring $D$: If $Z$ is a Zariski closed subset of the space of valuation rings dominating $D$, then $Z$ is connected if and only if the intersection of valuation rings in $Z$ is local.  The added strength here lies in the assertion that if the intersection $R$ of rings in $Z$ is local, the closed set $Z$ is connected. In particular, $Z$ need not be the entire set of valuation rings of $F$ that dominate $R$.  
%An example given in Section~{\bf [ref]} shows this stronger version of the connectedness theorem fails in dimension three.    

%We also prove versions of some of these  theorems for projective schemes rather than for valuation rings. 

\begin{notation}
We use the following notation.
\begin{enumerate} 
\item
If $R$ is a local ring, then ${\mathfrak{m}}_R$ denotes its maximal ideal.
\item
If $D$ is a subring of a field $F$, then 
$\Zar_F(D)$ is the set of valuation rings of $F$ that contain $D$.
\item

The set $\Dom_F(D)$ is the set of valuation rings $V$ in $\Zar_F(D)$ such that $\mathfrak m_V \cap D$ is a maximal ideal of $D$, i.e.\ $V$ dominates a localization of $D$ at a maximal ideal.  
%\item
% If $F$ is the quotient field of $D$, we write $\Zar(D)$ for $\Zar_F(D)$ and $\Dom_F(D)$ for $\Dom_F(D)$.  
 \item
If $X$ is a collection of local subrings of $F$, then $A(X)$ is the intersection of these rings, and $J(X)$ is the intersection of the maximal ideals of these rings. 
\item
If $X$ is a projective model of $F/D$  and $Z$ is a subset of $\Zar_F(D)$, then $X(Z)$ denotes the local rings in $X$ that are dominated by a valuation ring in $Z$.
\end{enumerate}

\end{notation}

   

\section{Projective models and the Zariski-Riemann space}



 In this section, we review the topology of the Zariski-Riemann space of a field and its relationship to projective models of the field.
Rather than use the language of projective schemes, we use the language of projective models to make the interplay between projective schemes and the Zariski-Riemann space of valuation rings easier to work with.

Let $F$ be a field,  let $D$ be a subring of $F$, and let Loc$_F(D)$ be the set of local subrings of $F$ that contain $D$ as a subring and have $F$ as their quotient field. The {\it Zariski topology} on Loc$_F(D)$ has a basis of  open sets the sets of the form 
    $${\mathcal{U}}(x_1,\ldots,x_n):=\{R\in {\rm Loc}(F/D):x_1,\ldots,x_n \in R\},$$
where $x_1,\ldots,x_n\in F$.  
The Zariski topology on Loc$_F(D)$ is quasicompact and the basic open sets  ${\mathcal{U}}(x_1,\ldots,x_n)$ are also quasicompact. In fact, by  \cite[Corollary~2.14]{MR3565806} or \cite[Example~2.1(7)]{MR3565813}, Loc$_F(D)$ is  a {\it spectral space},    a quasicompact $T_0$ topological space such that every open set is a union of  quasicompact open sets; the intersection of two quasicompact open sets is a quasicompact open set; and every irreducible closed subset has a unique generic point.  
(By a theorem of Hochster \cite{MR251026}, spectral spaces are   the topological spaces that are homeomorphic to the prime spectrum of a commutative ring.)  %(For the definition of a spectral space, see the appendix.)


The topology of a spectral space  can be refined to a Hausdorff topology called the {\it patch topology}. This topology has as a basis of open  sets the set of the form ${\mathcal{U}} \cap {\mathcal{V}}$, where ${\mathcal{U}}$ is a quasicompact open set  and ${\mathcal{V}}$ is the complement of a quasicompact open set. The sets ${\mathcal{U}} \cap {\mathcal{V}}$ are both closed and open, and it follows that the patch topology on $X$ is quasicompact, Hausdorff and zero-dimensional. 
%The {\it patch topology} on ${\rm Loc}(F/D)$ is the topology having as basis of open sets the quasicompact open sets  and their complements.
%The Zariski topology on  Loc$_F(D)$ is   Hausdorff only in the very special case in which $D$ is a field and $F$ is algebraic over $D$,  a situation we will not need to consider. However, the patch topology on Loc$_F(D)$ is Hausdorff in all cases and has a basis of clopen sets.
Closure of a set in a spectral space can be viewed as a generic closure of a patch closure. To state this, we recall that the {\it specialization order} on $X$ is given by $x \leq y$ if and only if $y$ is in the closure of $\{x\}$ with respect to the spectral topology on $X$.  
The next lemma is well known; a proof can be found in 
\cite[
%\href{https://stacks.math.columbia.edu/tag/0903}
{Tag 0903}]{stacks-project}.

\begin{lemma}  \label{generic closure} Let $Y$ be patch closed subspace of a spectral space $X$. The closure of a subspace $Y$ of $X$ in the spectral topology is the set of all $x \in X$ such that $y \leq x$ for some $y \in Y$. %
%
%is the set of local rings $R$ in $X$ such that $R$ is contained in a local ring in the patch closure of  $Y$. 
\end{lemma}
 


If $X$ is a subset of Loc($F/D$), 
 then the map $f:X \rightarrow \Spec  D $ that sends a local ring $R$ in $X$ to the prime ideal ${\mathfrak{m}}_R \cap D$ of $D$ is a \emph{spectral map}, i.e.\ continuous with respect to both the Zariski and patch topologies.
 In particular, if $X = \{D_P:P \in \Spec D\}$, then the map $X \rightarrow \Spec D $ is a homeomorphism, and so the Zariski topology on Loc$_F(D)$  can be viewed as extended from the Zariski topology on $\Spec D$.  
 
 
 

Patch closures of subspaces of Loc$_F(D)$ can be intricate to calculate; see for example 
\cite[Section~2]{MR3565813}.  
We work with two classes of subspaces of Loc$_F(D)$, those consisting of valuation rings and those consisting of local rings of points in projective schemes. 
To define the former, 
let  $\Zar_F(D)$ denote the set of valuation rings having quotient field $F$ and containing $D$ as a subring. The {\it Zariski-Riemann space} of $F/D$ is the set $\Zar_F(D)$ endowed with the subspace topology of the Zariski topology on Loc$_F(D)$. With this topology, $\Zar_F(D)$ is a spectral space. This 
%can be seen directly (e.g., as in \cite[Theorem 10.5, p.~74]{MR1011461}), and it 
follows   from the fact that $\Zar_F(D)$ is a patch closed subspace  of Loc$_F(D$); see \cite[Example 2.2(8)]{MR3565813} and \cite{MR0389876}. 

The other class of subspaces of Loc$_F(D)$ that we consider are those arising from  local rings of points in a projective integral scheme over $\Spec D $ with function field contained in $F$. To simplify the presentation, we view such schemes as projective models of $F/D$. Let $x_1,x_2,\ldots,x_n$ be nonzero elements of $ F$, and for each $i \in \{1,2,\ldots,n\}$, let $D_i = D[x_1x_i^{-1},\ldots,x_nx_i^{-1}]$.  Following \cite[p.~119]{MR0389876}, the {\it projective model of $F/D$ defined by $x_1,\ldots,x_n$} is the subspace of Loc$_F(D)$ with respect to the Zariski topology  consisting of all the localizations of the $D_i$, $i \in \{1,2,\ldots,n\}$, at the prime ideals of~$D_i$. 
%When $D$ has quotient field $F$ and $x_1,\ldots,x_n \in D$, then the projective model of $F/D$ defined by $x_1,\ldots,x_n$ is the blow-up Proj$(D[It])$ of $\Spec D$ along the subscheme defined by the ideal $I=(x_1,\ldots,x_n)D$. However, we will consider cases in which $D$ does not have quotient field $F$ and in this case  a projective model is not a blow-up of $\Spec D$. 

Since we will be passing back and forth between projective models and the Zariski-Riemann space, we will typically denote elements of projective models with Greek letters to help distinguish them. 
If $X$ is a  projective model  of $F/D$, then each valuation ring $V$ in $\Zar_F(D)$ {dominates} a unique local ring $\alpha$ in $X$. This property is referred to as {\it irredundance} in \cite[p.~115]{MR0389876}. Also, each local ring $\alpha$ in $X$ is dominated by a valuation ring in $\Zar_F(D)$, a property that is called {\it completeness} in  \cite[p.~127]{MR0389876}. (Irredundance and completeness reflect   the fact that a projective morphism is separated and proper.) 

%Thus there is a surjective mapping $d_X:\Zar_F(D)\rightarrow X$ that sends a valuation ring $V$ in $\Zar_F(D)$ to the local ring in $X$ that  $V$ dominates. 

\begin{lemma} \label{ccs}
{\em \cite[Lemma 4, p.~117]{MR0389876}} For each projective model $X$ of $F/D$, the   mapping $d_X:\Zar_F(D)\rightarrow X$ that sends a valuation ring $V$  in $\Zar_F(D)$ to the unique local ring in $X$ that  $V$ dominates  is surjective  and it is  continuous and closed   with respect to both the Zariski and patch topologies. 
\end{lemma}

The last assertion that $d_X$ is continuous and closed in the patch topology follows from \cite[Theorem~4.1]{MR932760} or \cite[Remark~3.2]{MR3299704}.
We denote by $X(Z)$ the image of $Z$ in $X$ under the mapping $d_X$.  
If $X$ and $Y$ are projective models of $F/D$, then $Y$ {\it dominates} $X$ if each local ring $\alpha$ in $Y$ dominates a local ring $\beta$ in $X$.  


A projective model $X$ can be identified with a locally ringed space (in fact, a projective scheme over $\Spec D$) in a natural way by equipping $X$ with a sheaf ${\mathcal{O}}_X$ such that for each nonempty Zariski open subset ${\mathcal{U}}$ of $X$, ${\mathcal{O}}_X({\mathcal{U}})$ is the intersection of the rings in ${\mathcal{U}}$. 
With this observation the next lemma follows from \cite[Exercise~II.3.2, p. 91]{MR0463157} and
\cite[Lemma~5, p.~119]{MR0389876}. 

\begin{lemma}  If $X$ and $Y$ are projective models of $F/D$ and $Y$ dominates $X$, then the map $d_{Y,X}:Y \rightarrow X$ that sends a local ring $\alpha$ in $Y$ to the unique local ring in $X$ that $\alpha$ dominates is surjective, closed and continuous in the Zariski and patch topologies. Moreover, $d_{Y,X} \circ d_Y = d_X$. 
\end{lemma} 

When equipped with the sheaf whose ring of sections over a  nonempty Zariski open subset  is the intersection of the rings in the  set,  $\Zar_F(D)$ is a locally ringed space and is the inverse limit of projective models. The following lemma is implicit in \cite{MR0389876}. For details, see \cite[Proposition~3.3(4)]{MR3299704}.

\begin{lemma} \label{inverse limit}   The set of projective models $Y$ of $F/D$ that dominate $X$ forms an inverse  system when equipped with the maps $d_{Y,X}$. The inverse limit of this system in the category of locally ringed spaces is $\Zar_F(D)$.  
\end{lemma} 

%In particular, each valuation ring $V \in $ $\Zar_F(D)$ is a directed union of the local rings in the projective models that are dominated by $V$.


%and if $Z$  is a closed subset of $\Zar_F(D)$, then $Z$ is the directed union of the sets $d_X(Z)$  \


% Examples of spectral spaces relevant to this article include Loc$_F(D)$, where $D$ is a subring of $F$; $\Zar_F(D)$; projective models of $F/D$; and $\Dom_F(D)$ when $D$ is local. For background and references for the assertion that the first three of these are spectral see {\bf [ref]}. That $\Dom_F(D)$ is spectral follows from the fact that $\Dom_F(D)$ is a closed set of the spectral space $\Zar_F(D)$. (It is closed since it is the preimage of a closed point under the continuous map $\Zar_F(D) \rightarrow \Dom_F(D)$.) 

%The topology of a spectral space $X$ can be refined to a Hausdorff topology called the {\it patch topology}. This topology has as a basis of open  sets the set of the form ${\mathcal{U}} \cap {\mathcal{V}}$, where ${\mathcal{U}}$ is a quasicompact open set  and ${\mathcal{V}}$ is the complement of a quasicompact open set. The sets ${\mathcal{U}} \cap {\mathcal{V}}$ are clopen (i.e., closed and open), and it follows that the patch topology on $X$ is quasicompact, Hausdorff and zero-dimensional. 

%A {\it patch closed} subset of a spectral space $X$ is  a subset of $X$ that is closed in the patch topology. 

%In light of Lemma~\ref{inverse limit}, it is useful to know when connectedness of subspaces of $\Zar_F(D)$ is equivalent to connectedness of the projections of these subspaces to projective models. To close this section, we give a topological criterion for this equivalence. This is done in Lemma~\ref{spectral connected 1}, and makes uses of two other lemmas that we record here because they will be needed later.  



%Some of our techniques   involve viewing the Zariski-Riemann space of a field as an inverse limit of projective models of the field. For this, we will use the following lemmas about inverse limits of spectral spaces. 


\section{Valuative connectedness}

In this section we prove one direction of  the valuative connectedness theorem stated in the introduction by proving that if $R$ is a local ring between the local ring $D$ and the field $F$ and $R$ dominates $D$ and is residually algebraic over $D$, then $\Dom_F(R)$ is a closed connected subset of $\Dom_F(D)$.  
This will be done by passing from $\Dom_F(D)$ to the projective models of $F/D$, where we can use Stein factorization to prove connectedness of fibers. 
For this reason, most of the results in this section deal with passing back and forth from $\Dom_F(D)$ to  relevant projective models of $F/D$. 
But first we make explicit a connection between residually algebraic local rings $R$ dominating $D$ and closed sets of $\Dom_F(D)$. Throughout the paper, working with closed connected sets rather than just connected sets is important for our methods. (This situation can always be arranged by replacing $D$ with a local ring between $D$ and $R$, over which $R$ is residually algebraic.)  



\begin{lemma} \label{J lemma} Let $D$ be a subring of the field $F$,  let $Z$ be a nonempty subset of $\Zar_F(D)$, and let $\overline{Z}$ be its Zariski closure in $\Zar_F(D)$.   Then $J(Z) = J(\overline{Z})$. 
\end{lemma}

\begin{proof}
 The patch closure $Z'$ of $Z$ is contained in $\overline{Z}$. Also, $\bigcap_{V \in Z}{\mathfrak{m}}_V = \bigcap_{V \in Z'}{\mathfrak{m}}_V$ by \cite[Lemma 2.3(1)]{MR3686983}. By Lemma~\ref{generic closure}, each valuation ring $V$ in $\overline{Z}$ is contained in a valuation ring $W$ in $Z'$. Since ${\mathfrak{m}}_W \subseteq  {\mathfrak{m}}_V$, it follows that $\bigcap_{V \in Z}{\mathfrak{m}}_V = \bigcap_{V \in \overline{Z}}{\mathfrak{m}}_V$. 
\end{proof}

\begin{proposition} \label{RA lemma} Let $D$ and $R$ be local subrings of
a field $F$ such that $R$  dominates $D$. 
Then $R$ is residually algebraic over $D$ if and only if $\Dom_F(R)$ is closed in $\Dom_F(D)$.
%
%Then $R = A(Z)$, where $Z$ is the Zariski closure of $\Dom_F(R)$ in $\Dom_F(D)$. 
\end{proposition} 

\begin{proof} 
 Suppose $R$ is residually algebraic over $D$. Then $R$ is integral over $D+{\mathfrak{m}}_R$, and so every valuation ring of $F$ containing $D+{\mathfrak{m}}_R$ contains $R$. Since $\Dom_F(R)$ is patch closed in $\Dom_F(D)$, the Zariski closure of $\Dom_F(R)$ in $\Dom_F(D)$ is the set of valuation rings $V \in \Dom_F(D)$ for which $V \subseteq U$ for some $U \in \Dom_F(R)$. Let $V$ be  such a valuation ring. 
% Let $U \in \Dom_F(R)$ and $V \in \Dom_F(D)$ with $V \subseteq U$. %Then $V$ is in the Zariski closure of $\Dom_F(R)$ in $\Dom_F(D)$. 
 By Lemma~\ref{J lemma},
 ${\mathfrak{m}}_R \subseteq J(\Dom_F(R)) \subseteq {\mathfrak{m}}_V$, and  so   
$V$ dominates $R$, which shows $V \in \Dom_F(R)$ and hence $\Dom_F(R)$ is Zariski closed in $\Dom_F(D)$.   

Conversely, if $\Dom_F(R)$ is Zariski closed in $\Dom_F(D)$, there exists a valuation ring $V \in \Dom_F(R)$ that is minimal over $D$, and hence is residually algebraic over $D$. Since $V$ dominates $R$, $R$ is residually algebraic over $D$. 
 \end{proof} 
 

The next lemma will be used to reduce arguments to the case of projective models. 

\begin{lemma} \label{directed models} \label{inverse system}
Let $D$ be a subring of the field $F$,  and let $Z$ be a nonempty subset of $\Zar_F(D)$.   
\begin{enumerate}
\item[$(1)$] The ring $A(Z)$ is the directed union of the rings $A(X(Z)),$ where $X$ ranges over the projective models of $F/D$.
%$$A(Z) =\bigcup_XA(X(Z)).$$

% where  $$A(Z) =\bigcup_XA(X(Z)),$$ 
 
\item[$(2)$] The  ring $A({Z})$ is local if and only if $A(X({Z}))$ is local for each projective model $X$ of $F/D$.
 

\item[$(3)$] If $Z = \bigcap_{i}Z_i$ for    an inverse system $\{Z_i\}$ of   patch closed subsets of $\Zar_F(D)$, then $\textstyle{A(Z)}$ is the directed union of the $ A(Z_i).$
\end{enumerate}
\end{lemma}

\begin{proof} 
Let $x_1,\ldots,x_n \in A(Z)$, and let $X$ be the projective model defined by the elements $1,x_1,\ldots,x_n$. Then $D[x_1,\ldots,x_n] \subseteq A(Z)$ and the localizations of $D[x_1,\ldots,x_n]$ at prime ideals are in $X$. 
Let $\alpha$ be a local ring in $X(Z)$.  Then $\alpha$ is dominated by a valuation ring $V$ in $Z$.  Since $D[x_1,\ldots,x_n] \subseteq V$, if ${\mathfrak{p}}$ is the center of $V$ in $D[x_1,\ldots,x_n]$, then $D[x_1,\ldots,x_n]_{\mathfrak{p}} \in X(Z)$. By the irredundance property discussed in Section~2, $\alpha = D[x_1,\ldots,x_n]_{\mathfrak{p}}$ since both of these local rings are in $X$ and are dominated by $V$. Thus $x_1,\ldots,x_n \in \alpha$, and since the choice of $\alpha$ in $X(Z)$ was arbitrary,  
  it follows that $x_1,\ldots,x_n \in A(X(Z))$. 
This proves that $A(Z)$ is contained in  the union of the $A(X(Z))$ and this union is directed. To prove the reverse inclusion, 
suppose $x \in A(X(Z))$ for some projective model $X$ of $F/D$. The fact that  $Z \subseteq d_X^{-1}(X(Z))$ implies
$x \in A(d_X^{-1}(X(Z))) \subseteq A(Z)$, and so $A(X(Z)) \subseteq A(Z)$ for all projective models $X$ of $F/D$. 

 

To see that (2) holds, note first that if  $A(X({Z}))$ is local for each projective model $X$ of $F/D$, then $A({Z})$ is a local ring since by (1), $A({Z})$ is the directed union of the rings $A(X({Z}))$, where $X$ ranges over the projective models of $F/D$.
Conversely, suppose $A(Z)$ is local, and let $X$ be a projective model of $F/D$. 
%For each $V \in {\mathcal{V}}$, let $R_V$ denote the ring in $X$ dominated by $V$. 
Let $a$ and $b$ be nonunits in $A(X({Z}))$. Since $A(X(Z))$ is an intersection of the  local rings in $X(Z)$, 
there are local rings $\alpha$ and $\beta$ in $X({Z})$ such that $a \in {\mathfrak{m}}_\alpha$ and $b \in {\mathfrak{m}}_\beta$. Since $\alpha,\beta \in X({Z})$, there are 
 valuation rings $V,W \in {Z}$ such that $V$ dominates $\alpha$ and $W$ dominates $\beta$.  Thus 
 $a \in {\mathfrak{m}}_V$ and $b \in {\mathfrak{m}}_W$, and  $a$ and $b$ are nonunits in $A(Z)$.  This implies  $a+b$ is a nonunit in the local ring $A(Z)$ and hence in the subring $A(X({Z}))$ of $A(Z)$, which shows the sum of nonunits in  $A(X({Z}))$ is itself not a unit. Therefore, $A(X({Z}))$ is local.  

 

For the proof of (3),  suppose $Z = \bigcap_i Z_i.$
 Since $Z \subseteq Z_i$ for each $i$, it is clear that $\bigcup_{i}A(Z_i)
\subseteq  A(Z)$. Suppose $x_1,\ldots,x_n \in A(Z)$, and let 
\begin{center}
${\mathcal{W}} = \{V \in \Zar_F(D):x_i \not \in V$ for some $1 \leq i \leq n\}$. 
\end{center}
Then $Z \cap {\mathcal{W}} = \emptyset$ since $x_1,\ldots,x_n \in A(Z)$, and so $\bigcap_{i} (Z_i \cap {\mathcal{W}}) = \emptyset$. 
Now $\Zar_F(D)$ is compact in the patch topology and  
the sets $Z_i \cap {\mathcal{W}}$ are patch closed subsets of  $\Zar_F(D)$, so there exist $i_1,\ldots,i_n$ such that $ Z_{i_1} \cap \cdots \cap Z_{i_n} \cap {\mathcal{W}} =\emptyset$. Choose $i$ such that  $Z_i \subseteq Z_{i_1} \cap \cdots \cap Z_{i_n}$. Then $Z_i \cap {\mathcal{W}} = \emptyset$, which implies $x_1,\ldots,x_n \in A(Z_i)$, proving that   $A(Z)$ is a directed union of the $A(Z_i).$
\end{proof}

If a  subspace  of a topological space  is connected, then so is its closure in the space. The next lemma shows that the converse is true in $\Zar_F(D)$. 

 
\begin{lemma} \label{dense} Let $D$ be a subring of a field $F$, and let $Z$ be a patch closed subset of $\Zar_F(D)$. Then $Z$ is connected if and only if the Zariski closure of $Z$ in $\Zar_F(D)$ is connected. 
%
%let $Y \subseteq Z$ be patch closed subspaces of $\Zar_F(D)$ such that $Y$ is Zariski dense in $ and let $Y$ be a Zariski dense subspace of $Z$ that is patch closed in $\Zar_F(D)$.  Then
%$Y$ is connected if and only if $Z$ is connected. 
\end{lemma}

\begin{proof}
The closure of a connected subspace of a topological space is always connected, so we need only show the inverse: if $Z$ is not connected, then neither is its Zariski-closure $\overline{Z}$.
Write $Z$ as the union of disjoint nonempty Zariski-closed subspaces $F_1, F_2$ of $Z$.
Since $\overline{Z}$ is the union of the nonempty Zariski-closed sets $\overline{F_1}, \overline{F_2}$, it suffices to show that $\overline{F_1}, \overline{F_2}$ are disjoint.

Suppose not, so there exists $y \in \overline{F_1} \cap \overline{F_2}$.
Since $Z$ is patch-closed set and $F_1, F_2$ are Zariski-closed subspaces, the sets $F_1, F_2$ are also patch-closed.
By Lemma~\ref{generic closure}, there exists $x_1 \in F_1, x_2 \in F_2$ such that $x_1 \le y$ and $x_2 \le y$.
Since the overrings of a valuation ring are totally ordered by inclusion, either $x_1 \le x_2$ or $x_2 \le x_1$, so assume without loss of generality that $x_1 \le x_2$.
But then, since $x_1 \in F_1$ and $F_1$ is a Zariski-closed subspace of $Z$, it follows that $x_2 \in F_1$, contradicting the fact that $F_1, F_2$ are disjoint.
\end{proof} 


% Since $\overline{Z}$ is connected, there is $V \in \overline{Z_1} \cap \overline{Z_2}$. Since $Z_1$ and $Z_2$ are Zariski closed in $Z$ and $Z$ is patch closed in $\Zar_F(D)$,  $Z_1$ and $Z_2$ are also patch closed in $\Zar_F(D)$.   By Lemma~\ref{generic closure}, $V \subseteq U_1 \cap U_2$ for some $U_1 \in Z_1$ and $U_2 \in Z_2$. 
%The overrings of a valuation ring are totally ordered under inclusion, so  
%  without loss of generality,
% $V \subseteq U_1 \subseteq U_2$. Since $U_1 \in Z$, $U_2 \in Z_2$ and $Z_2$ is Zariski closed in $Z$, we have $U_1 \in Z_2$. Thus $U_1 \in Z_1 \cap Z_2$, proving that $Z$ is not the disjoint union of nonempty Zariski closed subsets and hence that $Z$ is   connected.   

%  %The next  lemma is used in Section~3 to reduce the question of connectedness of a space of valuation rings to that of connectedness of fibers in  projective models. 

In general, an inverse limit of connected sets need not be connected. However, connectedness does lift to the inverse limit under the  hypotheses that we work with in this article. We give the topological argument for this in the next lemma, which  is surely known but for which we could not find a reference.
 
\begin{lemma} \label{spectral connected 1}
Let $S$ be the inverse limit of topological spaces $X_i$ such that the projection maps $d_i : S \rightarrow X_i$ are closed maps, and let $Z \subseteq S$ be a nonempty closed quasicompact subset.
Then $Z$ is connected if and only if $d_i (Z)$ is connected for each~$i$. 
\end{lemma}
 
\begin{proof}
Since the continuous image of a connected set is connected, if $Z$ is connected, then so is $d_i (Z)$ for each $i$.
Conversely, assume that $d_i ({Z})$ is connected for each $i$.
To show $Z$ is connected, for any pair of nonempty closed subsets $F_1, F_2 \subseteq Z$ such that $Z = F_1 \cup F_2$, we claim that $F_1 \cap F_2 \ne \emptyset$.
For each $i$, since $d_i$ is a closed map, we have that $d_i (F_1), d_i (F_2)$ are closed sets in $X_i$.
Since $d_i ({Z})$ is connected and $d_i ({Z}) = d_i (F_1) \cup d_i (F_2)$, the intersection $d_i (F_1) \cap d_i (F_2)$ must be nonempty.
For each $i$, define $Z_i$ to be the preimage of this intersection in $Z$, that is
    $$Z_i = Z \cap  d_i^{-1} (d_i (F_1) \cap d_i (F_2)).$$
Notice that for any $i, j$, if $i \ge j$, then $Z_i \subseteq Z_j$.
To see this, let $d_{j, i} : X_i \rightarrow X_j$ denote the compatibility map, where $d_{j, i} \circ d_i = d_j$.
Since $d_i (Z_i) \subseteq d_i (F_1) \cap d_i (F_2)$, by applying the map $d_{j, i}$, we obtain $d_j (Z_i) \subseteq d_j (F_1) \cap d_j (F_2)$.
Thus,
    $$Z_i \subseteq Z \cap d^{-1}_j (d_j (F_1) \cap d_j (F_2)) = Z_j.$$
This implies that the collection $\{ Z_i \}$ has the finite intersection property.
Indeed, let $j_1, \ldots, j_n$ be in the index set, and choose $j$ with $j \ge j_k$ for each $k$, $1 \le n$.
Then $Z_j$, which is nonempty, is a subset of each $Z_{j_k}$, and thus a subset of their intersection.

Since $\{ Z_i \}$ is a set of closed subsets of $Z$ with the finite intersection property and $Z$ is quasicompact, the intersection $\bigcap_{i} Z_i$ is nonempty. 
From the fact that $F_1, F_2$ are closed it follows\footnote{The general fact being used here is:
{\it
    Let $S$ be a topological space  that is an inverse limit of topological spaces $X_i$, and for each $i$ let $d_i: S \rightarrow X_i$ denote the projection map.
    If $Z$ is a closed subset of $S$, then $Z = \bigcap_{i} d_i^{-1}(d_i(Z))$.
}
To see that this is the case, observe that since the sets of the form $d_i^{-1} (Y)$, where $Y$ is a closed set in $X_i$, comprise a basis for the closed sets of $S$, it follows that $Z$ is the intersection of sets of the form $d_i^{-1}(Y)$, where $d_i(Z) \subseteq Y$ and $Y$ is closed in $X_i$. 
For such a set $Y$ we have $Z \subseteq d^{-1}_i(d_i(Z)) \subseteq d_i^{-1}(Y)$, and so $Z$ is the intersection of the sets of the form $d_i^{-1}(d_i(Z))$. 
 }
that $F_k =\bigcap_{i} d_i^{-1} (d_i (F_k))$ for $k = 1, 2$.
Thus 
\begin{eqnarray*} 
    \bigcap_i Z_i 
        &= & Z \cap \left( \bigcap_i d_i^{-1}(d_i(Z_1) \cap d_i(Z_2))\right) \\
        &= & Z \cap \left( \bigcap_i  d_i^{-1}(d_i(Z_1))\right) \cap \left( \bigcap_i d_i^{-1}( d_i(Z_2))\right) \\ 
        &= & Z_1 \cap Z_2,
\end{eqnarray*}
proving that $Z_1 \cap Z_2 \ne \emptyset$.
\end{proof}
 
%The lemma gives a way to  relate connectedness in $\Zar_F(D)$ with connectedness in the projective models of $F/D$:
 
\begin{lemma} \label{connected iff} Let $D$ be a subring of the field $F$.  A nonempty patch closed  subset ${Z}$ of $\Zar_F(D)$ is Zariski connected   if and only if $X({Z})$ is Zariski connected  for each projective model $X$ of $F/D$.  
\end{lemma} 

\begin{proof} 
If $Z$ is connected, then as a continuous image of $Z$, $X(Z)$ is  connected. Conversely, suppose $X(Z)$ is connected for all projective models $X$ of $F/D$.  Since the domination map $d_X$ is closed, the Zariski closure of $X(Z)$ in $X$ is $X(\overline{Z})$, where $\overline{Z}$ is the Zariski closure of $Z$ in $\Zar_F(D)$. Since a space containing a dense connected subspace is connected, $X(\overline{Z})$ is connected for all projective models $X$ of $F/D$.  
By Lemma~\ref{inverse limit}, the set of projective models of $F/D$ is an inverse system with respect to the domination maps for these models, and the inverse limit of this system is $\Zar_F(D)$. The domination maps $d_i:\Zar_F(D) \rightarrow X_i$ are the projection maps for the inverse limit, and by Lemma~\ref{ccs} these continuous maps are closed. Thus Lemma~\ref{spectral connected 1} implies $\overline{Z}$ is connected. By Lemma~\ref{dense}, $Z$ is connected. 
 \end{proof}

%\begin{example} Lemma~\ref{connected iff} is not valid if the assumption that $Z$ is closed is omitted. For example, suppose $D = k[[X,Y]]$, where $k$ is a field. Let $V$ be the order valuation ring of $D$ and let $U$ be a localization of $D$ at a height-one prime ideal of $D$. The set $X$ consisting of all the localizations of $D$ at a prime ideal is a projective model of $F/D$, where $F$ is the quotient field of $D$. Let $Z =\{V,U\}$. Since the valuation rings $V$ and $U$ are incomparable, $Z$ is not connected. However, $X(Z) = \{D,U\}$ is connected since $D \subseteq U$.   
%\end{example}
 


 
 
 

%If $Z$ is inverse closed, so is $\overline{Z}$.  Since Zariski closed sets are patch closed, $\overline{Z}$ consists of the valuation rings $V \in \Zar_F(D)$ such that $V \subseteq U$ for some $U \in Z$.  If $V \in \overline{Z}$ and $V \subseteq U$ for some $U \in \Dom_F(D)$, then there is $W \in Z$ such that $V \subseteq W$. If $U \subseteq W$, then since $W \in \overline{Z}$, $U \in \overline{Z}$.   Otherwise, if $W \subseteq U$, since $W \in Z$ and $Z$ is inverse closed, $U \in Z$.  




%For an integral domain $D$, denote by $\mathfrak{X}(D)$ the set of valuation rings $V$ in $\Zar(D)$ that are centered on maximal ideals in $D$.  Then $\mathfrak{X}(D)$ is a Zariski closed subspace of $\Zar(D)$ since it is patch closed in $\Zar(D)$ and a down set in $\Zar(D)$.   

%\begin{lemma} Let $R$ be a local subring of a field $F$, let $K$ be a subfield of $F$ and $D$ a local subring of $K$ dominated by $R$.  The map $d:\Zar_F(R) \rightarrow \Zar(K/D)$ is surjective and  continuous and closed in the Zariski topology.
%\end{lemma} 






 
%  \begin{remark}
%\label{Noetherian local remark}
%Theorem~\ref{Noetherian local} implies Zariski's connectedness theorem since if $D$ is as in the corollary and $X$ is a projective model of $F/D$, then the fiber of the  map $X \rightarrow \Spec D$ over  the closed point of $\Spec D$ is the image of the map $\Dom_F(D) \rightarrow X$, so that since $\Dom_F(D)$ is connected, so is this fiber. 
%\end{remark} 



% \begin{lemma} \label{same res} Let $D$ be a local subring of a field $F$,  let $R$ be an integrally closed local subring of $F$ that dominates $D$, and let $Z$ be the Zariski closure of $ \Dom_F(R)$ in $\Dom_F(D)$. If $D$ and $R$ have the same residue field, then $R = A(\overline{Z})$. 
% \end{lemma}
 
% \begin{proof} 
% Since $R$ is local and has the same residue field as $D$, we have $R = D+{\mathfrak{m}}_R$. Since $R$ is integrally closed, ${\mathfrak{m}}_R = J(\Dom_F(R))$.   
% By Lemma~\ref{J lemma}, $J(
% \Dom_F(R)) = J({Z})$, and so $R =  D+{\mathfrak{m}}_R = D+J({Z}) \subseteq A({Z}) \subseteq A(\Dom_F(R)) = R$, proving that $R = A({Z})$. 
% \end{proof}





Using the previous lemmas, we now     prove that local rings give rise to connected sets by reducing to   the case of projective models, where Stein factorization is available. 

\begin{theorem} \label{first direction} Suppose $D \subseteq R$ are local subrings of a field $F$ such that $R$ dominates $D$, $R$ is integrally closed in $F$ and $R$ is residually algebraic over $D$. Then   $\Dom_F(R)$ is a closed connected subset   of $\Dom_F(D)$.  
\end{theorem}

\begin{proof} 


Let $Z=\Dom_F(R)$. To show $Z$ is connected, it suffices by  Lemma~\ref{connected iff} to show $X(Z )$ is connected for each projective model $X$ of $F/R$. Let $X$ be a projective model of $F/R$. 
As follows from the property of completeness discussed in Section~2 (or from the valuative criterion for properness), the image $X(Z)$ of $Z$ in $X$ under the domination map $d_X:\Zar_F(D) \rightarrow X$ is the fiber of
 the map $X \rightarrow \Spec R$ over the closed point 
  of $\Spec R$. Since  $R$ is integrally  closed in $F$, 
$R$ is integrally closed in the   
   function field of $X$. 
Now we apply the non-Noetherian version of Stein factorization from \cite[{Tag~03H2}]{stacks-project}, which in this case implies that since $R$ is  local and integrally closed in $F$, the fiber $X(Z)$ of the map
  $X \rightarrow \Spec R$ over the closed point 
  of $\Spec R$ is
     connected in the Zariski topology.
%  Since $d_X$ is a closed map,   the image $X( {Z})$ of $Z$ under $d_X$ is the Zariski closure of  $X(Z)$ in $X$.  
This proves
    ${Z}$ is connected. That $Z$ is closed follows from Proposition~\ref{RA lemma}. 
%Apply Lemmas~\ref{Noetherian local} and~\ref{RA lemma} and use the fact that a set is connected if it has a dense connected subset. 
\end{proof} 
 



\section{Local intersections over closed and connected sets} 

The purpose of this section is to prove a   converse to  Zariski's  connectedness theorem (Theorem~\ref{proj case}) and to lift this converse to spaces of valuation rings (Theorem~\ref{is local}).  

%If $X$ is a projective model over $D$, then the {\it exceptional locus} of $X$ is the (closed) set of local rings in $X$ that dominate $D$. 


\begin{theorem} \label{proj case} 
Let $D$ be a local subring of a field $F$ such that $\Spec D$ is a Noetherian space, and let $X$ be a projective model of $F/D$.
\begin {enumerate}
\item[$(1)$]
If $Y$ is a Zariski closed subset of the set of local rings in $X$ that dominate $D$, then $A(Y)$ is a semilocal ring whose maximal ideals are contracted from the maximal ideals of the rings in $Y$.  
\item[$(2)$]
If in addition  $Y$ is connected,  then $A({Y})$ is a   local   ring that is dominated by each ring in $Y$. %If also the model $X$ is normal, then $A$ is a normal Noetherian ring. 
\end{enumerate}
\end{theorem}

\begin{proof} 
Since $\Spec D$ is Noetherian, each finitely generated $D$-algebra has  Noetherian prime spectrum \cite[Section~2]{MR229627}, and it follows that $X$, as a finite union of Noetherian spaces, is a Noetherian space. Thus $Y$, as a closed subset of $X$, is the union of finitely many closed irreducible subsets ${Y}_1,{Y}_2,\ldots,{Y}_n$ of $X$. %If ${Y}$ consists of a single local ring, then clearly $A({Y})$ is a local ring, so we assume ${Y}$ has more than one member. 
%Since ${Y}$ is connected, each irreducible component of ${Y}$ has infinitely many members, and so, as a closed set, ${Y}$ contains local rings that are not closed points in $X$. 
For each $i$, let $U_i$ be the local ring that is the unique generic point of ${Y}_i$.  Then $U_i$ is 
essentially of finite type over $D$, and  $U_i$  dominates $D$. Also, ${Y}_i = \{\alpha \in {Y}:\alpha \subseteq U_i\}$. 
%
%and are residually transcendental over $D$.  
The Zariski closed set ${Y}$ thus consists of the rings in $X$ that are contained in at least one of the $U_i$. 

%Since $X$ is a nonsingular projective model over $D$ and the local rings in ${Y}$ dominate $D$, there are at most two one-dimensional  local rings in ${Y}$ {\bf [ref]}. It follows that ${Y}$ is either the exceptional fiber of $X$ over $D$ or ${Y}$ is an irreducible component whose generic point is a prime divisor of $D$. 

%In the former case, $A({Y})$ is the intersection of all the two-dimensional regular local rings in $X$ and hence is $D$ itself. Thus we need only consider the latter case in which ${Y}$ is an irreducible component whose generic point is a prime divisor $U$ of $D$.
 
 Let $A = A({Y})$.
Then $A \subseteq U_1 \cap \cdots \cap U_n$. For each $i \in \{1,2,\ldots,n\}$,  let $P_i = {\mathfrak{m}}_{U_i} \cap A$. We claim each $P_i$ is a maximal ideal of $A$. 
%and 
%let ${Y}_i = \{\alpha \in {Y}:\alpha \subseteq U_i\}$ be the irreducible component of ${Y}$ whose generic point is $U_i$. 
Observe that for each $i$, $$
D / \mathfrak{m}_D =
(D+{\mathfrak{m}}_{U_i})/{\mathfrak{m}}_{U_i}
\subseteq 
(A+{\mathfrak{m}}_{U_i})/{\mathfrak{m}}_{U_i}
\subseteq \bigcap_{\alpha \in {Y}_i}(\alpha+ {\mathfrak{m}}_{U_i})/{\mathfrak{m}}_{U_i} \subseteq U_i/{\mathfrak{m}}_{U_i}.$$
Since $U_i$ dominates $D$, the maximal ideal ${\mathfrak{m}}_{U_i}$ of $U_i$ contracts in $D$ to the maximal ideal of $D$. 
By \cite[(6.2.19),p.~161]{MR1617523}, the collection $\{(\alpha+{\mathfrak{m}}_{U_i})/{\mathfrak{m}}_{U_i} \mid \alpha \in {Y}_i\}$ is a projective model over the residue field of $D$. Since the intersection of local rings in a projective model is contained in the intersection of all the valuation rings of the function field of the model, the extension     $$(D+{\mathfrak{m}}_{U_i})/{\mathfrak{m}}_{U_i} \subseteq
\bigcap_{\alpha \in {Y}_i}(\alpha+ {\mathfrak{m}}_{U_i})/{\mathfrak{m}}_{U_i}$$ 
is an integral extension whose base ring is a field.
This implies the intermediate ring $(A+{\mathfrak{m}}_{U_i})/{\mathfrak{m}}_{U_i}$ 
  is a field. Since $P_i = A \cap {\mathfrak{m}}_{U_i}$, we conclude $P_i$ is a maximal ideal of~$A$.  
 

%There is $i \in \{1,2,\ldots,n\}$ such that $\alpha \subseteq U_i$, so ${\mathfrak{m}}_{U_i} \cap \alpha \subseteq {\mathfrak{m}}_{\alpha}$. Thus $P_i = {\mathfrak{m}}_{U_i} \cap A  \subseteq {\mathfrak{m}}_{\alpha} \cap A$. Since $P_i$ is maximal in $A$, we have $x,y \in {\mathfrak{m}}_\alpha \cap A = P_i$, contrary to the chose

Next we claim that the $P_i$ are all the maximal ideals of $A$. 
Let $x \in A$ be a nonunit. Since $A = A (Y)$, there exists $\alpha \in Y$ such that $x \in \mathfrak{m}_{\alpha}$.
Since $\alpha \subseteq U_i$ for some $i$ and  $\alpha \cap \mathfrak{m}_{U_i} \subseteq \mathfrak{m}_{\alpha}$,
it follows that $ P_i = A \cap \mathfrak{m}_{U_i} \subseteq A \cap \mathfrak{m}_{\alpha}$.
Since $P_i$ is a maximal ideal of $A$, we conclude that  $P_i = A \cap \mathfrak{m}_{\alpha}$, and so $x \in P_i$. This proves that every nonunit in $A$ is contained in one of the $P_i$. Using prime avoidance, we conclude that each proper ideal in $A$ is contained in one of the $P_i$, and hence the $P_i$ are all the maximal ideals of $A$.  This proves item 1. 


 

  %Next we claim $P_1$ is the unique maximal ideal of $A$. Suppose by way of contradiction that  there is $p \in P_1$ such that $1+p$ is not a unit in $A$. Then since $A = A({Y})$ there is $\alpha \in {Y}$ such that $1+p$ is in ${\mathfrak{m}}_\alpha$. Since $\alpha \subseteq U_i$ for some $i$, we have $\alpha \cap  {\mathfrak{m}}_{U_i} \subseteq {\mathfrak{m}}_\alpha$, and hence $P_1=P_i = A \cap {\mathfrak{m}}_{U_i} \subseteq A \cap {\mathfrak{m}}_\alpha.$ Since $P_1$ is a maximal ideal of $A$, it follows that $1+p \in A \cap {\mathfrak{m}}_\alpha = P_1$, a contradiction to the assumption that $p \in P_1$. Therefore, $P_1$ is the unique maximal ideal of $A$.  
  %The ring $A$ is residually algebraic over $D$ since it is dominated by local rings that are residually algebraic over $D$.  It is  two-dimensional since it is dominated by two-dimensional local rings that have prime ideals that do not dominate $D$. 
  
For item 2,   suppose $Y$ is connected. To prove that $A(Y)$ is local, we show that $P_1 = P_2 = \cdots = P_n$.
Suppose this is not the case. Since the $P_i$ are maximal ideals we may assume without loss of generality that $P_1 \ne P_2 \cap \cdots \cap P_n$.
Since each $P_i$ is maximal, there exist nonzero elements $x \in P_1$ and $y \in P_2 \cap \cdots \cap P_n$ such that $x + y = 1$.
%Since $xy \in P_1 \cap \cdots P_n$, 
For an element $t\in F$, let $${\mathcal{V}}(t) =\{ \alpha \in X \mid t \notin \alpha \},$$ and note that ${\mathcal{V}}(t)$ is a Zariski closed subset of $X$.
Now $U_1 \in {\mathcal{V}}(x^{-1})$ and $U_2,\ldots,U_n \in {\mathcal{V}}(y^{-1})$, so 
since $U_1,\ldots,U_n$ are the generic points for the irreducible closed components of ${Y}$, we have  ${Y} \subseteq {\mathcal{V}}(x^{-1}) \cup {{\mathcal{V}}}(y^{-1})$. 
%Since $x+y=1$, we have that $x \not \in P_2 \cap \cdots \cap P_n$ and $y \not \in P_1$.  Without loss of generality, assume $x \not \in P_2$. Then $x$ is a unit in $U_2$ but not $U_1$ and $y$ is a unit in $U_1$.  
To see that $Y \cap \mathcal{V} (x^{-1}) \cap \mathcal{V} (y^{-1}) = \emptyset$, suppose $\alpha \in {Y} \cap  {\mathcal{V}}(x^{-1}) \cap {\mathcal{V}}(y^{-1})$.
Since $x, y \in A \subseteq \alpha$ and $x^{-1}, y^{-1} \notin \alpha$, both $x$ and $y$ are nonunits in $\alpha$ with $x + y = 1$, contradicting the locality of~$\alpha$.
From this we conclude ${Y}$ is the disjoint union of the nonempty closed subsets $ {\mathcal{V}}(x^{-1}) \cap {Y}$ and ${\mathcal{V}}(y^{-1}) \cap {Y}$, a contradiction to the fact that  ${Y}$ is connected. Therefore, $P_1 = P_2 = \cdots = P_n$, which proves $A$ is local. Moreover, $A$ is dominated by each $\alpha \in {Y}$ since the previous argument shows  $P_1  = A 
\cap {\mathfrak{m}}_\alpha$.
\end{proof} 




We illustrate Theorem~\ref{proj case} with an example from an interesting  geometric setting. The details for the example can be found in \cite[Example~7.6]{MR4097907}.  

  
\begin{example} Let $D$ be a two-dimensional regular local ring with quotient field $F$, and let ${\mathfrak m} =(x,y)D$ be the maximal ideal of $D$.  The ring  $R = D[y^2/x]_{(x,y,y^2/x)D[y^2/x]}$ is a two-dimensional Noetherian normal local domain that is not regular, since its maximal ideal ${\mathfrak{m}}_R=(x,y,y^2/x)R$ requires three generators. The blow-up  Proj $R[{\mathfrak m}_Rt]$ of ${\mathfrak m}_R$ is the desingularization of $R$, which we may view as a  projective model of $F/R$. It is a subset of the  projective model $X= $ Proj $D[Jt]$, where $J = (x^2, xy, y^3)D$. Let $Y$ be the set of local rings in $X$ that dominate $R$. Then $Y$ is a closed subset of $X$, and it is connected since it is irreducible. 
Moreover, $R = A(Y)$, as in 
 Theorem~\ref{proj case}.
\end{example}

More generally, if $D$ is a local Noetherian domain and $J$ is an ideal with blowup  $X = $ {Proj~$D[Jt]$}, then for each closed connected subset $Y$ of the fiber in $X$ over the maximal ideal of $D$, the local rings in $Y$ 
 intersect to a local ring, according to Theorem~\ref{proj case}. In particular, the irreducible closed sets in $Y$ yield local rings. 


We   lift  Theorem~\ref{proj case} to the Zariski-Riemann space of valuation rings by using the fact from Lemma~\ref{inverse limit} that this locally ringed space is the inverse limit   of projective models. 




\begin{theorem} \label{is local}
Let $D$ be a local subring of the field $F$.
If ${Z}$ is a nonempty Zariski closed and connected subspace of  $\Dom_F(D)$, then $A({Z})$ is a local domain that dominates $D$ and is residually algebraic over $D$. 
\end{theorem}

\begin{proof}
Let $D'$ be the prime subring of $D$. 
We may write $D $ as a directed union $D= \bigcup_{i} D_i$, where for each $i$, $D_i$ is a local ring essentially of finite type over $D'$ and dominated by $D$.  
Then $\Dom_F(D) = \bigcap_{i} \Dom_F(D_i)$.
 For each $i$, let 
 $Z_i$ be the Zariski closure of ${Z}$ in  $\Dom_F(D_i)$. Then $Z_i \cap \Dom_F(D) = Z$ since $Z$ is Zariski closed in $\Dom_F(D)$, so 
 %By Lemma~\ref{closed}, each valuation ring in $Z_i$ dominates $D_i$.
  $Z = \bigcap_i Z_i$. Since $\{Z_i\}$ is an inverse system of Zariski closed (hence patch closed) subsets of $\Zar_F(D)$,
  Lemma~\ref{inverse system}(3) implies $A(Z)$ is a directed union of the $A(Z_i)$. Thus by Lemma~\ref{directed models}(1), $A(Z)$ is the directed union of the rings $A(X(Z_i))$, where $X$ ranges over the  projective models of $F/D_i$ and $i$ ranges over the index set of $\{D_i\}$. 
  
  To prove that $A(Z)$ is local it suffices to prove that $A({X}(Z_i))$ is local for each~$i$ and  each projective model $X$.  
  Let $i$ be in the index set, and let $X$ be a projective model of $F/D_i$. 
 Since $Z_i$ contains the dense connected set $Z$, $Z_i$ is connected. 
 For each   projective model $X$ of $F/D_i$, the map $d_X:\Zar_F(D_i) \rightarrow X$  that sends a valuation ring to the unique local ring it dominates in $X$ is by Lemma~\ref{ccs}  continuous, closed and surjective.
 Thus, since $Z_i$ is    closed and connected, so is the subset $d_{X}(Z_i)=X(Z_i)$ of the set of local rings in $X$ dominating $D_i$,
 and hence by Theorem~\ref{proj case},  $A(X(Z_i))$ is a local ring. Therefore, $A(Z)$ is a local ring.  

 Finally, since each valuation ring in $Z$ dominates $D$, we have ${\mathfrak{m}}_D \subseteq J(Z) \subseteq {\mathfrak{m}}_{A(Z)}$, and so $A(Z)$  dominates $D$.  
 Let $U$ be a valuation ring in $Z$. Since $Z$ is closed in $\Dom_F(D)$, there is a valuation ring $V \subseteq U$ such that $V \in Z$, $V$ is minimal over $D$ and hence $V$ is residually algebraic over $D$.  Thus since $D +{\mathfrak{m}}_V \subseteq A(Z)+{\mathfrak{m}}_V  \subseteq V$, it follows that $A(Z)$ is residually algebraic over $D$. 
\end{proof}


Combining Theorems~\ref{first direction} and~\ref{is local}, we obtain the valuative connectedness theorem from the introduction. 


  
 
 
 Theorem~\ref{is local}, and hence the valuative connectedness theorem,  are not true if the assumption that $Z$ is Zariski closed is removed.  Example~\ref{Nagata example} shows that even  patch closed in place of Zariski closed  is not strong enough for the theorem to remain valid. 


\begin{example} \label{Nagata example} Let $D$ be a Noetherian normal local domain with Krull dimension at least $2$ and quotient field $F$. Let $x$ be an indeterminate for $F$, and for each subring $R$ of $F$, let $R(x)$ denote the Nagata function ring of $R$, i.e.\ $R(x)$ is the set of rational functions $f(x)/g(x)$ in $F(x)$, where $f(x),g(x) \in R[x]$ and the ideal generated by the coefficients of $g$ is $R$. Let $Z = \Dom_F(D)$. By Theorem~\ref{first direction}, $Z$ is connected in the Zariski topology.  Let $Z^* = \{V(x):V \in Z\}$. Then $Z^*$ is homeomorphic to $Z$ (cf.~\cite[Proposition~4.2(b)]{MR3299704}) and patch closed in $\Dom_{F(x)}(D(x))$ \cite[Proposition~5.6(6)]{MR3299704}, yet $A(Z^*)$ is a B\'ezout  domain whose maximal ideals are the centers of the minimal valuation rings in $Z^*$ \cite[Corollary~5.8]{MR3299704}. 
Since $D$ is Noetherian and  has dimension at least $2$, there are infinitely many minimal valuation rings in $Z$, hence in $Z^*$, and so $A(Z^*)$ has infinitely many maximal ideals.  
Thus $Z^*$ is a patch closed, Zariski connected subset of $\Dom_{F(x)}(D(x))$ for which $A(Z^*)$ is not local. 
\end{example} 

As the example shows,  Theorem~\ref{is local} can fail if $Z$ is not Zariski closed. However,  the theorem can still be applied  
to the Zariski closure of $Z$ and information about $A(Z)$  obtained.  



 \begin{corollary} \label{is local corollary} Let $D$ be a local subring of the field $F$ such that $D$ is not a field. If $Z$ is a nonempty connected subspace of $\Dom_F(D)$, then
 there is a local integrally closed subring of $A(Z)$ whose complete integral closure contains $A(Z)$.  

 \end{corollary} 
 
 \begin{proof} 
 Let $\overline{Z}$ denote the Zariski closure of $Z$ in $\Dom_F(D)$.
 Then $\overline{Z}$ is connected since it has a dense connected subspace, and so $A(\overline{Z})$ is by Theorem~\ref{is local} an integrally closed local domain.
 By Lemma~\ref{J lemma}, $J(Z) = J(\overline{Z})$. 
 Since the valuation rings in $Z$ dominate $D$ and $D$ is not a field, there is $0 \ne x \in \bigcap_{V \in Z}{\mathfrak{m}}_V$. Consequently, $xA(Z) \subseteq A(\overline{Z})$, so that $A(Z)$ is in the complete integral closure of $A(\overline{Z})$. \end{proof}
 
 
 
%\begin{corollary}
%Let $D$ be a subring of a field $F$, and let $Z$ be a subset of $\Zar_F(D)$ such that there is a valuation ring $V$ in $Z$ having the same residue field as $D$. If $Z$ is connected, then $A(Z)$ is local. 




 
%\begin{remark} Since an analytically unramified   Noetherian local domain $D$ with quotient field $F$ has the property that the integral closure of each finitely generated $D$-subalgebra    of $F$ is a finitely generated $D$-module, and hence a Noetherian ring {\bf [ref]},  a variation on the proof of Corollary~\ref{Nagata ring} shows that if 
% $R$ is an integrally closed local overring of $D$ having the same residue field as $D$,  then $\Dom_F(R)$  is connected. 
  
%\end{remark}



 




%The sets $\bigcap_{\alpha}Y(\V_\alpha) \setminus Y(\V)$ are patch closed and since $d_Y$ is a spectral map, the preimage of this set is patch closed also. By compactness there are projective models $Y_1,\ldots,Y_n$ such that $d^{-1}_{Y_i}( (\bigcap_\alpha Y_i(\V_\alpha) \setminus Y_i(\V)) = \emptyset$. 

% Let $Y$ be the join of the $Y_i$. 
 
% Then $d^{-1}_{Y}( (\bigcap_\alpha Y(\V_\alpha) \setminus Y(\V)) = \emptyset$, and so $\bigcap_\alpha Y(\V_\alpha) = Y(\V).$
 
% Since $D$ is Noetherian, $X$ is a Noetherian space, and hence satisfies the descending chain condition on closed sets. Consequently, the intersection of all the $X(\V_\alpha)$ is equal to $X(\V_{\alpha_1}) \cap \cdots \cap X(\V_{\alpha_n})$ for some $\alpha_1,\ldots,\alpha_n$. Since the collection $\{D_\alpha\}$ is directed, there is $\beta$ such that $D_{\alpha_1},\ldots,D_{\alpha_n} \subseteq D_\beta$, and hence $\V_\beta \subseteq  \V_{\alpha_1} \cap \cdots \cap \V_{\alpha_n},$ which in turn implies $X(V_\beta) \subseteq  X(\V_{\alpha_1}) \cap \cdots \cap X(\V_{\alpha_n}).$ Therefore, $X(V_\beta)$ is contained in $X(V_\alpha)$ for all $\alpha$.




 



 
 %\begin{corollary} Let $D$ be a Noetherian local domain. A completely integrally closed overring $A$ of $D$ is local if and only if there is a connected subspace $\V$ of $\Dom_F(D)$ such that $A = A(\V)$. 
 %\end{corollary} 
 
% \begin{proof} 
% Suppose $A$ is local,  let $\V = \Dom(A)$ and let $\overline{\V}$ denote the Zariski closure of $\V$ in $\Dom_F(D)$. Then $A(\overline{\V}) \subseteq A(\V) = A$. Observe that $J:=\bigcap_{V \in \V}{\mathfrak{m}_V} \subseteq A$.   Since the rings in $\V$ dominate $D$, there exists $0 \ne x \in J$. Now $xA(\V) \subseteq A$
 
% Since $\V$ is connected 
 
 %\end{proof}
 
%\todo[inline]{Can we remove ``Noetherian'' from the corollary?}



%\begin{appendices}

%\section{Spectral spaces} 




 
 
%\end{appendices}

%\nocite{*}
\printbibliography



\end{document}


\begin{proof} If ${\mathcal{V}}$ is connected, then the closure $\overline{{\mathcal{V}}}^\alpha$ of ${\mathcal{V}}$ in $\Zar_F(D_\alpha)$ is connected for each $\alpha$  since a dense subset of a connected set is connected. 

Conversely, suppose $\Dom_F(D_\alpha)$ is connected for each $\alpha$.   To show  ${\mathcal{V}}$ is connected,
suppose by way of contradiction that ${\mathcal{V}}$ is a disjoint union of two nonempty closed sets $ {\mathcal{V}}_1 $ and ${\mathcal{V}}_2$. For each $\alpha$, since $\overline{{\mathcal{V}}}^\alpha$ is connected and $\overline{{\mathcal{V}}}^\alpha = \overline{{\mathcal{V}}_1}^\alpha \cup \overline{{\mathcal{V}}_2}^\alpha$, 
 it follows that $Z_{\alpha}:=\overline{Z_1}^\alpha \cap \overline{Z_2}^\alpha$ is nonempty. 
 

Conversely, suppose  $\overline{{\mathcal{V}}}^\alpha$ is connected for each $\alpha$.   
To show  ${\mathcal{V}}$ is connected,
suppose by way of contradiction that ${\mathcal{V}}$ is a disjoint union of two nonempty closed sets $ {\mathcal{V}}_1 $ and ${\mathcal{V}}_2$.  For each $\alpha$, since $\overline{{\mathcal{V}}}^\alpha$ is connected and $\overline{{\mathcal{V}}}^\alpha = \overline{{\mathcal{V}}_1}^\alpha \cup \overline{{\mathcal{V}}_2}^\alpha$, 
 it follows that $Z_{\alpha}:=\overline{Z_1}^\alpha \cap \overline{Z_2}^\alpha$ is nonempty. 
 
If $\alpha_1,\ldots,\alpha_n$ are in the index set for $\{D_\alpha\}$, then the intersection of the 
$Z_{\alpha_i}$ is nonempty since if  
 $\alpha$ is chosen such that $D_{\alpha_1},\ldots,D_{\alpha_n} \subseteq D_\alpha$,  
then $Z_\alpha \subseteq 
Z_{\alpha_1} \cap \cdots \cap V_{\alpha_n}$. 
Thus   the collection   $\{{\mathcal{V}}_\alpha\}$ has the finite intersection property. 

has the finite intersection property. 
Since ${\mathcal{V}}$ is quasicompact (it is a closed subset of the quasicompact space Zar($F/D$)), there is a valuation ring $V \in \bigcap_{X}{\mathcal{V}}_X$.
Since ${\mathcal{V}}$ is a disjoint union of ${\mathcal{V}}_1$ and $ {\mathcal{V}}_2$, we can 
assume without loss of generality that
 $V \not \in  {\mathcal{V}}_2$. 
 
 For each projective model $X$ of $F/D$, since $d_X(V) \in X({\mathcal{V}}_1) \cap X({\mathcal{V}}_2)$, there is a valuation ring $V_X \in {\mathcal{V}}_2$ such that $d_X(V_X) = d_X(V)$. We claim $V$ is a Zariski limit point of the $V_X$. Let ${\mathcal{U}}$ be an open neighborhood of $V$.  There are $x_1,\ldots,x_m \in F$ such that $V \in  {\mathcal{U}}(x_1,\ldots,x_m) \subseteq {\mathcal{U}}.$ There is a projective model $X$ of $F/D$ such that the localizations of $D[x_1,\ldots,x_n]$ at prime ideals are all in $X$.  Since $d_X(V)$ is a localization of $D[x_1,\ldots,x_n]$ at a prime ideal and  $d_X(V) = d_X(V_X)$, we have $D[x_1,\ldots,x_n] \subseteq V_X$, so that $V_X \in {\mathcal{U}}(x_1,\ldots,x_m)$. 
 This shows
 $V$ is a limit point of the $V_X$, where $X$ ranges over the projective models of $F/D$.   But
 the $V_X$ are all in the closed set ${\mathcal{V}}_2$ while the limit point $V$ is not. This contradiction shows ${\mathcal{V}}$ is connected. 
\end{proof} 


************


\begin{lemma} \label{proj case} 
Let $D$ be a local subring of a field $F$ such that $\Spec D$ is a Noetherian space, and let $X$ be a projective model of $F/D$. If ${Y}$ is    a  closed connected subset  of the set of local rings in $X$ that dominate $D$, then $A({Y})$ is a   local   ring that is dominated by each ring in $Y$. %If also the model $X$ is normal, then $A$ is a normal Noetherian ring. 
\end{lemma}

\begin{proof} 
Since $\Spec D$ is Noetherian, each finitely generated $D$-algebra has a Noetherian prime spectrum {\bf [Ohm-Pendleton]}, and it follows that $X$, as a finite union of Noetherian spaces, is a Noetherian space. Thus $Y$, as a closed subset of $X$, is the union of finitely many closed irreducible subsets ${Y}_1,{Y}_2,\ldots,{Y}_n$ of $X$. %If ${Y}$ consists of a single local ring, then clearly $A({Y})$ is a local ring, so we assume ${Y}$ has more than one member. 
%Since ${Y}$ is connected, each irreducible component of ${Y}$ has infinitely many members, and so, as a closed set, ${Y}$ contains local rings that are not closed points in $X$. 
For each $i$ let $U_i$ be the local ring that is the unique generic point of ${Y}_i$.  Then $U_i$ is 
essentially of finite type over $D$ and dominates $D$. Also, ${Y}_i = \{\alpha \in {Y}:\alpha \subseteq U_i\}$. 
%
%and are residually transcendental over $D$.  
The Zariski closed set ${Y}$ thus consists of the rings in $X$ that are contained in at least one of the $U_i$. 

%Since $X$ is a nonsingular projective model over $D$ and the local rings in ${Y}$ dominate $D$, there are at most two one-dimensional  local rings in ${Y}$ {\bf [ref]}. It follows that ${Y}$ is either the exceptional fiber of $X$ over $D$ or ${Y}$ is an irreducible component whose generic point is a prime divisor of $D$. 

%In the former case, $A({Y})$ is the intersection of all the two-dimensional regular local rings in $X$ and hence is $D$ itself. Thus we need only consider the latter case in which ${Y}$ is an irreducible component whose generic point is a prime divisor $U$ of $D$.
 
 Let $A = A({Y})$.
Then $A \subseteq U_1 \cap \cdots \cap U_n$. For each $i \in \{1,2,\ldots,n\}$,  let $P_i = {\mathfrak{m}}_{U_i} \cap A$. We claim each $P_i$ is a maximal ideal of $A$. 
%and 
%let ${Y}_i = \{\alpha \in {Y}:\alpha \subseteq U_i\}$ be the irreducible component of ${Y}$ whose generic point is $U_i$. 
Observe that for each $i$, $$(D+{\mathfrak{m}}_{U_i})/{\mathfrak{m}}_{U_i}
\subseteq 
(A+{\mathfrak{m}}_{U_i})/{\mathfrak{m}}_{U_i}
\subseteq \bigcap_{\alpha \in {Y}_i}(\alpha+ {\mathfrak{m}}_{U_i})/{\mathfrak{m}}_{U_i} \subseteq U_i/{\mathfrak{m}}_{U_i}.$$
Since $U_i$ dominates $D$, the maximal ideal ${\mathfrak{m}}_{U_i}$ of $U_i$ contracts in $D$ to the maximal ideal of $D$. 
By {\bf [ref]}, the collection $\{(\alpha+{\mathfrak{m}}_{U_i})/{\mathfrak{m}}_{U_i} \mid \alpha \in {Y}_i\}$ is a projective model over the residue field $(D+{\mathfrak{m}}_{U_i})/{\mathfrak{m}}_{U_i}$ of $D$, so by {\bf [ref]} we have    $$(D+{\mathfrak{m}}_{U_i})/{\mathfrak{m}}_{U_i} = \bigcap_{\alpha \in {Y}_i}(\alpha+ {\mathfrak{m}}_{U_i})/{\mathfrak{m}}_{U_i},$$ which implies  $(D+{\mathfrak{m}}_{U_i})/{\mathfrak{m}}_{U_i}=(A+{\mathfrak{m}}_{U_i})/{\mathfrak{m}}_{U_i}$ 
 and  hence this ring is a field. Since $P_i = A \cap {\mathfrak{m}}_{U_i}$, we conclude $P_i$ is a maximal ideal of $A$.  
 
Next we show  $P_1 = P_2 = \cdots = P_n$.
Suppose not. Since the $P_i$ are maximal ideals we may assume without loss of generality that $P_1 \ne P_2 \cap \cdots \cap P_n$.
Since each $P_i$ is maximal, there exist nonzero elements $x \in P_1$ and $y \in P_2 \cap \cdots \cap P_n$ such that $x + y = 1$.
%Since $xy \in P_1 \cap \cdots P_n$, 
For an element $t\in F$, let $${\mathcal{V}}(t) =\{ R \in X \mid t \notin R \},$$ and note that ${\mathcal{V}}(t)$ is a closed subset of $X$.
Now $U_1 \in {\mathcal{V}}(x^{-1})$ and $U_2,\ldots,U_n \in {Y}(y^{-1})$, so 
since $U_1,\ldots,U_n$ are the generic points for the irreducible closed components of ${Y}$, we have  ${Y} \subseteq {Y}(x^{-1}) \cup {Y}(y^{-1})$. 
%Since $x+y=1$, we have that $x \not \in P_2 \cap \cdots \cap P_n$ and $y \not \in P_1$.  Without loss of generality, assume $x \not \in P_2$. Then $x$ is a unit in $U_2$ but not $U_1$ and $y$ is a unit in $U_1$.  
To see that $Y \cap \mathcal{V} (x^{-1}) \cap \mathcal{V} (y^{-1}) = \emptyset$, suppose $R \in {Y} \cap  {\mathcal{V}}(x^{-1}) \cap {\mathcal{V}}(y^{-1})$.
Since $x, y \in A \subseteq R$ and $x^{-1}, y^{-1} \notin R$, both $x$ and $y$ are nonunits in $R$ with $x + y = 1$, contradicting the locality of $R$.
From this we conclude ${Y}$ is the disjoint union of the nonempty closed subsets $ {\mathcal{V}}(x^{-1}) \cap {Y}$ and ${\mathcal{V}}(y^{-1}) \cap {Y}$, a contradiction to the fact that  ${Y}$ is connected. Therefore, $P_1 = P_2 = \cdots = P_n$.  



%There is $i \in \{1,2,\ldots,n\}$ such that $\alpha \subseteq U_i$, so ${\mathfrak{m}}_{U_i} \cap \alpha \subseteq {\mathfrak{m}}_{\alpha}$. Thus $P_i = {\mathfrak{m}}_{U_i} \cap A  \subseteq {\mathfrak{m}}_{\alpha} \cap A$. Since $P_i$ is maximal in $A$, we have $x,y \in {\mathfrak{m}}_\alpha \cap A = P_i$, contrary to the chose

Finally, we claim that $P_1$ is the unique maximal ideal of $A$.
Let $x \in A$ be a nonunit. Since $A = A (Y)$, there is $\alpha \in Y$ such that $x \in \mathfrak{m}_{\alpha}$.
Since $\alpha \subseteq U_i$ for some $i$ and  $\alpha \cap \mathfrak{m}_{U_i} \subseteq \mathfrak{m}_{\alpha}$,
it follows that $P_1 = P_i = A \cap \mathfrak{m}_{U_i} \subseteq A \cap \mathfrak{m}_{\alpha}$.
Since $P_1$ is a maximal ideal of $A$, we conclude $P_1 = A \cap \mathfrak{m}_{\alpha}$, so $x \in P_1$, which proves $P_1$ is the unique maximal ideal of $A$. 
Moreover, $A$ is dominated by each $\alpha \in {Y}$ since the previous argument shows  $P_1  = A 
\cap {\mathfrak{m}}_\alpha$.
  %Next we claim $P_1$ is the unique maximal ideal of $A$. Suppose by way of contradiction that  there is $p \in P_1$ such that $1+p$ is not a unit in $A$. Then since $A = A({Y})$ there is $\alpha \in {Y}$ such that $1+p$ is in ${\mathfrak{m}}_\alpha$. Since $\alpha \subseteq U_i$ for some $i$, we have $\alpha \cap  {\mathfrak{m}}_{U_i} \subseteq {\mathfrak{m}}_\alpha$, and hence $P_1=P_i = A \cap {\mathfrak{m}}_{U_i} \subseteq A \cap {\mathfrak{m}}_\alpha.$ Since $P_1$ is a maximal ideal of $A$, it follows that $1+p \in A \cap {\mathfrak{m}}_\alpha = P_1$, a contradiction to the assumption that $p \in P_1$. Therefore, $P_1$ is the unique maximal ideal of $A$.  
  %The ring $A$ is residually algebraic over $D$ since it is dominated by local rings that are residually algebraic over $D$.  It is  two-dimensional since it is dominated by two-dimensional local rings that have prime ideals that do not dominate $D$. 
\end{proof} 


\begin{theorem} \label{Dom(R)} 
Let $D$ be a normal Noetherian local domain, and let 
 $R$ be an integrally closed local domain with quotient field $F$ such that $R$  is the directed union of normal  Noetherian $D$-subalgebras that have quotient field $F$.  If  $\Dom_F(R)$ is closed in $\Zar(D)$, then $\Dom_F(R)$ is connected. 
\end{theorem}


\begin{proof} 
%Let $Z = \Dom_F(R)$, and let $\overline{Z}$  be the Zariski closure of $Z$ in $\Dom_F(D)$. Since $Z$ is patch closed,   Lemma~\ref{up} implies that to prove $Z$ is connected in the inverse topology, it suffices to prove $\overline{Z}$ is connected in the inverse topology. 
%
We may write $A(\overline{Z}) = \bigcup_i D_i$ for a directed set of normal Noetherian local $D$-subalgebras of $A(\overline{Z})$, each of which is dominated by $A(\overline{Z})$ and has quotient field $F$. 
For each $i$, let $Z_i=\Dom_F(D_i)$.  We first observe that each $Z_i$ is connected. 
For each~$i$ and
 projective model $X$ of $F/D_i$, {\bf [ref]} implies $X(Z_i)$ is the fiber of
 the map $X \rightarrow \Spec D_i$ over the closed point 
  of $\Spec D_i$. By Zariski's Connectedness Theorem, the fact that $D_i$ is normal, Noetherian and local implies  $X(Z_i)$ is   connected in the Zariski topology, and hence also by Lemma {\bf [ref]} connected in the inverse topology. Thus  $Z_i$ is connected in the inverse topology by Lemma~\ref{connected iff} {\bf [modify]}.

Let $Z = \Dom_F(R)$.  Since $Z$ is closed, 
 Lemma~\ref{connected iff} implies that to prove $Z$ is connected, it suffices to show $X(Z)$ is a connected set for each projective model $X$ of $F/D$.  
Let $X$ be a projective model of $F/D$.
Letting $Y$ range over the projective models that dominate $X$, we have by Lemma~\ref{closure} and the fact that $Z$ and the $Z_i$ are patch closed,  
\begin{center} 
$Z = \bigcap_Y d_Y^{-1}(Y(Z))$ \: and \: $Z_i =  \bigcap_Y d_Y^{-1}(Y(Z_i))$ for all $i$.\end{center}  
%Since the image of a connected set under a continuous map is connected, to show $X(Z)$ is connected it suffices to show 
% $X(Z) = X(Z_i)$ for some $i$.  
Also, since $R$ is the directed union of the $D_i$, we have $Z = \bigcap_{i}Z_i$. Therefore, again letting $Y$ vary over all projective models of $F/D$ that dominate $X$,  we have
\begin{eqnarray} \label{eq d}
\bigcap_{Y,i}d_Y^{-1}\left( Y(Z_i) \setminus Y(Z)\right)
%\bigcap_i \bigcap_{Y}d_Y^{-1}( %Y(Z_i) \setminus Y(Z)) 
& \subseteq & \bigcap_i \left(
\bigcap_{Y}d_Y^{-1}(Z_i)\right)\setminus \bigcap_{Y}d_Y^{-1}(Z) \\ 
&\subseteq & \left(\bigcap_i Z_i\right) \setminus Z = 
 \emptyset.
\end{eqnarray}
For each $i$, $Z_i$ is a patch closed subset of $\Zar_F(D)$, as is $Z$. For each projective model $Y$ of $F/D$, the map $d_Y$ is continuous  in the patch topology {\bf [ref]}, and since $Y$ is Hausdorff in the patch topology, this map is closed in the patch topology. Thus $Y(Z_i) = d_Y(Z_i)$ is patch closed, as is $Y(Z)$.  Since $Y$ is a Noetherian space, the Zariski closed sets in $Y$ are   patch clopen sets in $Y$. Since $d_Y$ is closed in the Zariski topology,  $Y(Z)$ is a patch clopen set in $Y$. Thus for each $i$, $Y(Z_{i}) \setminus Y(Z)$ is a patch closed subset of $Y$. Consequently, 
$d_Y^{-1}\left( Y(Z_i) \setminus Y(Z)\right)$ is a patch closed set in $\Zar_F(D)$. 

These considerations, along with the fact that $\Zar_F(D)$ is quasicompact in the patch topology {\bf [ref]}, imply via equation~(\ref{eq d}) that there are projective models $Y_1,\ldots,Y_m$ and indices $i_1,\ldots,i_n$  such that
$$\bigcap_{j,k}d_{Y_j}^{-1}\left( Y_j(Z_{i_k}) \setminus Y_j(Z)\right) = \emptyset.$$ Let $Y$ be the join of $Y_1,\ldots,Y_m$. Then 
\begin{center} $d_{Y}^{-1}\left(Y(Z_{i_k})\setminus Y(Z)\right) \: \subseteq \: d_{Y_j}^{-1}(Y_j(Z_{i_k})\setminus Y_j(Z))$ \:for all $j,k$,
\end{center}
and so 
\begin{center} $d_{Y}^{-1}(\bigcap_{k}Y(Z_{i_k})\setminus Y(Z)) = \emptyset$
\end{center}
Since $d_Y$ is surjective, this implies $\bigcap_{k}Y(Z_{i_k})= Y(Z).$ Since $\{D_i\}$ is a directed set, we may choose $i$ such that $D_{i_1},\ldots,D_{i_n} \subseteq D_i.$ For each $k$, $D_i$ dominates $D_{i_k}$ and so  $$Z_i = \Dom_F(D_i) \subseteq \Dom_F(D_{i_k}) = Z_{i_k}.$$ 
Thus $$d_Y^{-1}(Y(Z_i) \setminus Y(Z)) \: \subseteq \: \bigcap_{k}
d_Y^{-1}(Y(Z_{i_k}) \setminus Y(Z_{i})) = \emptyset.$$ It follows that $Y(Z_i) = Y(Z)$, and since $Y$ dominates $X$, we conclude $X(Z_i) = X(Z).$ We have established that  $Z_i$ is connected, and so  since continuous images of connected sets are connected, 
 this proves $X(Z_i) = X(Z)$ is connected. Therefore, $Z$ is connected by Lemma~\ref{connected iff}.
 \end{proof}
 
 
 
 \bigskip
 
 By {\bf [ref]}, the domination mapping $d_X$ is inversely continuous for each projective model $X$ of $F/D$. Thus if ${Z}$ is inversely connected, then $X({Z})$ is inversely connected for each projective model $X$ of $F/D$  since a continuous image of a connected set is connected. 
 
 Conversely, suppose  $X({Z})$ is inversely connected for each projective model $X$ of $F/D$. To show  ${Z}$ is inversely connected, 
suppose by way of contradiction that ${Z}$ is a disjoint union of two nonempty inversely closed sets $ {Z}_1 $ and ${Z}_2$.  For each projective model $X$ of $F/D$, 
$X(Z_1)$ and $X(Z_2)$ are inversely closed sets in $X(Z)$ since 



the domination mapping $d_X:\Zar_F(D) \rightarrow X$ is an inversely  closed map. 
Thus since $X({Z})$ is connected  and
$X({Z}) = X({Z}_1) \cup X({Z}_2)$,  it follows that $X({Z}_1) \cap X({Z}_2)$ is nonempty. Let $d_X:\Zar_F(D) \rightarrow X$ denote the domination map that sends a valuation ring in $\Zar_F(D)$ to its center. 
For each projective model $X$, define a closed subset $Z_X$ of $Z$ by    $${Z}_X:= Z \cap  d_X^{-1}(X({Z}_1) \cap X({Z}_2)).$$
The set $Z_X$ is nonempty since  $X({Z}_1) \cap X({Z}_2)$ is nonempty and $Z =Z_1 \cup Z_2$.  
%Then $  {Z}_X \subseteq  {Z}$ for all choices of $X$. 

We claim that  if $X_1,\ldots,X_n$ are projective models of $F/D$, then the intersection of the ${Z}_{X_i}$ is nonempty. Let $X$ be the join of the $X_i$, as in
\cite[p.~120]{MR0389876}. 
  Then $X$ dominates each $X_i$, and the image of $X({Z}_1)$ in $X_i$ under the domination map $d_{X,X_i}:X \rightarrow X_i$ is $X_i({Z}_1)$. Similarly, the image of $X({Z}_2)$ is $X_i({Z}_2)$. 
If $V \in {Z}_X$, then $V$ dominates a local ring in 
$X({Z}_1) \cap X({Z}_2)$, and hence also in $X_i({Z}_1) \cap X_i({Z}_2)$ for each $i$. 
Thus $V \in Z_{X_i}$ and it 
 follows that ${Z}_X \subseteq {Z}_{X_1} \cap \cdots \cap {Z}_{X_n}.$
Since  ${Z}_{X}$ is nonempty, we conclude ${Z}_{X_1} \cap \cdots \cap  {Z}_{X_n}$ is nonempty. This proves the collection   $\{{Z}_X:X$ is a projective model$\}$ of  closed subsets ${Z}_X$  has the finite intersection property. 
Since ${Z}$ is quasicompact (it is a closed subset of the quasicompact space Zar($F/D$)), there is a valuation ring $V \in \bigcap_{X}{Z}_X$.
Since ${Z}$ is a disjoint union of ${Z}_1$ and $ {Z}_2$, we can 
assume without loss of generality that
 $V \not \in  {Z}_2$. 
 
 For each projective model $X$ of $F/D$, since $d_X(V) \in X({Z}_1) \cap X({Z}_2)$, there is a valuation ring $V_X \in {Z}_2$ such that $d_X(V_X) = d_X(V)$. We claim $V$ is a Zariski limit point of the $V_X$. Let ${\mathcal{U}}$ be an open neighborhood of $V$.  There are $x_1,\ldots,x_m \in F$ such that $V \in  {\mathcal{U}}(x_1,\ldots,x_m) \subseteq {\mathcal{U}}.$ The projective model $X$ of $F/D$ defined by $1,x_1,\ldots,x_m$ has the property that the the localizations of $D[x_1,\ldots,x_n]$ at prime ideals are all in $X$.  Since $d_X(V)$ is a localization of $D[x_1,\ldots,x_n]$ at a prime ideal and  $d_X(V) = d_X(V_X)$, we have $D[x_1,\ldots,x_n] \subseteq V_X$, so that $V_X \in {\mathcal{U}}(x_1,\ldots,x_m)$. 
 This shows
 $V$ is a limit point of the $V_X$, where $X$ ranges over the projective models of $F/D$.   But
 the $V_X$ are all in the closed set ${Z}_2$ while the limit point $V$ is not. This contradiction shows ${Z}$ is connected. 
\end{proof} 


 \begin{corollary} \label{au} Let $D$ be an analytically unramified   Noetherian local  domain with quotient field $F$, and let $R$ be an integrally closed local  overring of $D$.
   If $\Dom_F(R)$ is closed in $\Zar(D)$, then $\Dom_F(R)$ is connected. 
 
 \end{corollary} 
 
 \begin{proof}  Write $R = \bigcup_i D_i$ for a directed set of Noetherian local $D$-subalgebras of $R$, each of which is dominated by $R$ and is essentially of finite type over $D$. Since $D$ is analytically unramified, the integral closure of $D_i$ is Noetherian for each $i$.  
Thus by replacing $D_i$ with the localization of its integral closure at a maximal ideal contained in the maximal ideal of $R$, we may assume the $D_i$ are normal Noetherian local domains. Now apply Theorem~\ref{Dom(R)}.
 \end{proof} 


\begin{proof}  We may write $R = \bigcup_i D_i$ for a directed set of normal Noetherian local $D$-subalgebras of $R$, each of which is dominated by $R$ and has quotient field $F$. 
For each $i$, let $Z_i=\Dom_F(D_i)$.  We first observe that each $Z_i$ is connected. 
For each $i$ and
 projective model $X$ of $F/D_i$, {\bf [ref]} implies $X(Z_i)$ is the fiber of
 the map $X \rightarrow \Spec D_i$ over the closed point 
  of $\Spec D_i$. By Zariski's Connectedness Theorem, the fact that $D_i$ is normal, Noetherian and local implies  $X(Z_i)$ is   connected. Thus  $Z_i$ is connected by Lemma~\ref{connected iff}.

Let $Z = \Dom_F(R)$.  Since $Z$ is closed, 
 Lemma~\ref{connected iff} implies that to prove $Z$ is connected, it suffices to show $X(Z)$ is a connected set for each projective model $X$ of $F/D$.  
Let $X$ be a projective model of $F/D$.
Letting $Y$ range over the projective models that dominate $X$, we have by Lemma~\ref{closure} and the fact that $Z$ and the $Z_i$ are patch closed,  
\begin{center} 
$Z = \bigcap_Y d_Y^{-1}(Y(Z))$ \: and \: $Z_i =  \bigcap_Y d_Y^{-1}(Y(Z_i))$ for all $i$.\end{center}  
%Since the image of a connected set under a continuous map is connected, to show $X(Z)$ is connected it suffices to show 
% $X(Z) = X(Z_i)$ for some $i$.  
Also, since $R$ is the directed union of the $D_i$, we have $Z = \bigcap_{i}Z_i$. Therefore, again letting $Y$ vary over all projective models of $F/D$ that dominate $X$,  we have
\begin{eqnarray} \label{eq d}
\bigcap_{Y,i}d_Y^{-1}\left( Y(Z_i) \setminus Y(Z)\right)
%\bigcap_i \bigcap_{Y}d_Y^{-1}( %Y(Z_i) \setminus Y(Z)) 
& \subseteq & \bigcap_i \left(
\bigcap_{Y}d_Y^{-1}(Z_i)\right)\setminus \bigcap_{Y}d_Y^{-1}(Z) \\ 
&\subseteq & \left(\bigcap_i Z_i\right) \setminus Z = 
 \emptyset.
\end{eqnarray}
For each $i$, $Z_i$ is a patch closed subset of $\Zar_F(D)$, as is $Z$. For each projective model $Y$ of $F/D$, the map $d_Y$ is continuous  in the patch topology {\bf [ref]}, and since $Y$ is Hausdorff in the patch topology, this map is closed in the patch topology. Thus $Y(Z_i) = d_Y(Z_i)$ is patch closed, as is $Y(Z)$.  Since $Y$ is a Noetherian space, the Zariski closed sets in $Y$ are   patch clopen sets in $Y$. Since $d_Y$ is closed in the Zariski topology,  $Y(Z)$ is a patch clopen set in $Y$. Thus for each $i$, $Y(Z_{i}) \setminus Y(Z)$ is a patch closed subset of $Y$. Consequently, 
$d_Y^{-1}\left( Y(Z_i) \setminus Y(Z)\right)$ is a patch closed set in $\Zar_F(D)$. 

These considerations, along with the fact that $\Zar_F(D)$ is quasicompact in the patch topology {\bf [ref]}, imply via equation~(\ref{eq d}) that there are projective models $Y_1,\ldots,Y_m$ and indices $i_1,\ldots,i_n$  such that
$$\bigcap_{j,k}d_{Y_j}^{-1}\left( Y_j(Z_{i_k}) \setminus Y_j(Z)\right) = \emptyset.$$ Let $Y$ be the join of $Y_1,\ldots,Y_m$. Then 
\begin{center} $d_{Y}^{-1}\left(Y(Z_{i_k})\setminus Y(Z)\right) \: \subseteq \: d_{Y_j}^{-1}(Y_j(Z_{i_k})\setminus Y_j(Z))$ \:for all $j,k$,
\end{center}
and so 
\begin{center} $d_{Y}^{-1}(\bigcap_{k}Y(Z_{i_k})\setminus Y(Z)) = \emptyset$
\end{center}
Since $d_Y$ is surjective, this implies $\bigcap_{k}Y(Z_{i_k})= Y(Z).$ Since $\{D_i\}$ is a directed set, we may choose $i$ such that $D_{i_1},\ldots,D_{i_n} \subseteq D_i.$ For each $k$, $D_i$ dominates $D_{i_k}$ and so  $$Z_i = \Dom_F(D_i) \subseteq \Dom_F(D_{i_k}) = Z_{i_k}.$$ 
Thus $$d_Y^{-1}(Y(Z_i) \setminus Y(Z)) \: \subseteq \: \bigcap_{k}
d_Y^{-1}(Y(Z_{i_k}) \setminus Y(Z_{i})) = \emptyset.$$ It follows that $Y(Z_i) = Y(Z)$, and since $Y$ dominates $X$, we conclude $X(Z_i) = X(Z).$ We have established that  $Z_i$ is connected, and so  since continuous images of connected sets are connected, 
 this proves $X(Z_i) = X(Z)$ is connected. Therefore, $Z$ is connected by Lemma~\ref{connected iff}.
 \end{proof}
 
 \begin{corollary} \label{au} Let $D$ be an analytically unramified   Noetherian local  domain with quotient field $F$, and let $R$ be an integrally closed local  overring of $D$.
   If $\Dom_F(R)$ is closed in $\Zar(D)$, then $\Dom_F(R)$ is connected. 
 
 \end{corollary} 
 
 \begin{proof}  Write $R = \bigcup_i D_i$ for a directed set of Noetherian local $D$-subalgebras of $R$, each of which is dominated by $R$ and is essentially of finite type over $D$. Since $D$ is analytically unramified, the integral closure of $D_i$ is Noetherian for each $i$.  
Thus by replacing $D_i$ with the localization of its integral closure at a maximal ideal contained in the maximal ideal of $R$, we may assume the $D_i$ are normal Noetherian local domains. Now apply Theorem~\ref{Dom(R)}.
 \end{proof} 
 
 
  \begin{corollary} 
 Let $D$ be Nagata  domain, and let $R$ be an integrally closed overring of $D$ such that each maximal ideal of $R$ contracts to a maximal ideal of $D$.  
   Then $R$ 
   is local if and only if $\Dom_F(R)$ is a   connected set. 
 \end{corollary}
 
 \begin{proof}
 If $R$ is local, then $\Dom_F(R)$ is a Zariski connected set by Theorem~\ref{Dom(R)}. For the converse, observe that $\Dom_F(D)$ is a Zariski closed subset of $\Zar(D)$ and $\Dom_F(R)$ is a Zariski closed subset of $\Dom_F(D)$, so that $\Dom_F(R)$ is a Zariski closed subset of $\Zar(D)$. Therefore, $R$ is local by  Theorem~\ref{is local}.
 \end{proof} 
 
 
 
  \begin{proof} 
 Since $\overline{Z}$ is closed, 
 Lemma~\ref{connected iff} implies that to prove $\overline{Z}$ is connected, it suffices to show $d_\alpha(\overline{Z})$ is a connected set for each $\alpha$. Since 
 
 
Let $X$ be a projective model of $F/D$.
Letting $Y$ range over the projective models that dominate $X$, we have by Lemma~\ref{closure} and the fact that $Z$ and the $Z_i$ are patch closed,  
\begin{center} 
$Z = \bigcap_Y d_Y^{-1}(Y(Z))$ \: and \: $Z_i =  \bigcap_Y d_Y^{-1}(Y(Z_i))$ for all $i$.\end{center}  
%Since the image of a connected set under a continuous map is connected, to show $X(Z)$ is connected it suffices to show 
% $X(Z) = X(Z_i)$ for some $i$.  
Also, since $R$ is the directed union of the $D_i$, we have $Z = \bigcap_{i}Z_i$. Therefore, again letting $Y$ vary over all projective models of $F/D$ that dominate $X$,  we have
\begin{eqnarray} \label{eq d}
\bigcap_{Y,i}d_Y^{-1}\left( Y(Z_i) \setminus Y(Z)\right)
%\bigcap_i \bigcap_{Y}d_Y^{-1}( %Y(Z_i) \setminus Y(Z)) 
& \subseteq & \bigcap_i \left(
\bigcap_{Y}d_Y^{-1}(Z_i)\right)\setminus \bigcap_{Y}d_Y^{-1}(Z) \\ 
&\subseteq & \left(\bigcap_i Z_i\right) \setminus Z = 
 \emptyset.
\end{eqnarray}
For each $i$, $Z_i$ is a patch closed subset of $\Zar_F(D)$, as is $Z$. For each projective model $Y$ of $F/D$, the map $d_Y$ is continuous  in the patch topology {\bf [ref]}, and since $Y$ is Hausdorff in the patch topology, this map is closed in the patch topology. Thus $Y(Z_i) = d_Y(Z_i)$ is patch closed, as is $Y(Z)$.  Since $Y$ is a Noetherian space, the Zariski closed sets in $Y$ are   patch clopen sets in $Y$. Since $d_Y$ is closed in the Zariski topology,  $Y(Z)$ is a patch clopen set in $Y$. Thus for each $i$, $Y(Z_{i}) \setminus Y(Z)$ is a patch closed subset of $Y$. Consequently, 
$d_Y^{-1}\left( Y(Z_i) \setminus Y(Z)\right)$ is a patch closed set in $\Zar_F(D)$. 

These considerations, along with the fact that $\Zar(D)$ is quasicompact in the patch topology {\bf [ref]}, imply via equation~(\ref{eq d}) that there are projective models $Y_1,\ldots,Y_m$ and indices $i_1,\ldots,i_n$  such that
$$\bigcap_{j,k}d_{Y_j}^{-1}\left( Y_j(Z_{i_k}) \setminus Y_j(Z)\right) = \emptyset.$$ Let $Y$ be the join of $Y_1,\ldots,Y_m$. Then 
\begin{center} $d_{Y}^{-1}\left(Y(Z_{i_k})\setminus Y(Z)\right) \: \subseteq \: d_{Y_j}^{-1}(Y_j(Z_{i_k})\setminus Y_j(Z))$ \:for all $j,k$,
\end{center}
and so 
\begin{center} $d_{Y}^{-1}(\bigcap_{k}Y(Z_{i_k})\setminus Y(Z)) = \emptyset$
\end{center}
Since $d_Y$ is surjective, this implies $\bigcap_{k}Y(Z_{i_k})= Y(Z).$ Since $\{D_i\}$ is a directed set, we may choose $i$ such that $D_{i_1},\ldots,D_{i_n} \subseteq D_i.$ For each $k$, $D_i$ dominates $D_{i_k}$ and so  $$Z_i = \Dom_F(D_i) \subseteq \Dom_F(D_{i_k}) = Z_{i_k}.$$ 
Thus $$d_Y^{-1}(Y(Z_i) \setminus Y(Z)) \: \subseteq \: \bigcap_{k}
d_Y^{-1}(Y(Z_{i_k}) \setminus Y(Z_{i})) = \emptyset.$$ It follows that $Y(Z_i) = Y(Z)$, and since $Y$ dominates $X$, we conclude $X(Z_i) = X(Z).$ Since $d_X$ is a closed map, the Zariski closure of $X(Z_i)$ is $X(\overline{Z_i})$ and the Zariski closure of $X(Z)$ is $X(\overline{Z})$.
We have established that  $Z_i$ is connected, and so  since continuous images of connected sets are connected,  $X(Z_i) = X(Z)$ is connected. 
Thus $X(\overline{Z})$ is connected, which proves  $\overline{Z}$ is connected by Lemma~\ref{connected iff}.
 \end{proof}
 
 
  \begin{lemma} \label{Noetherian space} Let
$d:S \rightarrow X$ be a closed spectral map of spectral spaces with $X$ a Noetherian space. If $\{Z_i\}$ is an inverse system  of patch closed sets of $S$, then $d(\bigcap_i{Z_i}) = d(Z_j)$ for some $j$. 
\end{lemma} 



\begin{proof} 
Let $Z = \bigcap_i \overline{Z_i}$.

Since $X$ is a Noetherian space and the $Z_i$ form an inverse system, it follows that there is $j$ such that $d({Z})
\subseteq d(\bigcap_i \overline{Z}_i)
\subseteq  \bigcap_i d(\overline{Z_i}) = d(\overline{Z_j}).$ Thus 


\bigskip

Let $Y = \bigcap_{i}d(Z_i)$.  Since $d$ is a closed map each $d(Z_i)$ is closed. Since $X$ is a Noetherian space, the closed sets satisfy the descending chain condition, and hence there are $i_1,\ldots,i_n$ such that $Y = d(Z_{i_1}) \cap \cdots \cap d(Z_{i_n})$. Since $\{Z_i\}$ is an inverse system of closed sets of $S$, there is $i$ with $Z_i \subseteq Z_{i_1} \cap \cdots \cap Z_{i_n}$. Thus $$Y \subseteq d(Z_i) \subseteq d(Z_{i_1} \cap \cdots \cap Z_{i_n}) \subseteq  d(Z_{i_1}) \cap \cdots \cap d(Z_{i_n}) = Y.$$ Hence $d(Z_i) = Y$, which proves the lemma.


 
%$Z =  d^{-1}(Y(Z))$ \: and \: $Z_i =  \bigcap_Y d_Y^{-1}(Y(Z_i))$ for all $i$.\end{center}  
%Since the image of a connected set under a continuous map is connected, to show $X(Z)$ is connected it suffices to show 
% $X(Z) = X(Z_i)$ for some $i$.  

 Let $Z = \bigcap_{i}Z_i$. 
 
 

\begin{eqnarray} \label{eq d}
\bigcap_{i}d^{-1}\left( d(Z_i) \setminus d(Z)\right)
%\bigcap_i \bigcap_{Y}d_Y^{-1}( %Y(Z_i) \setminus Y(Z)) 
& \subseteq & \bigcap_i \left(
d^{-1}(d(Z_i))\right)\setminus d^{-1}(d(Z)) \\ 
&\subseteq & \left(\bigcap_i Z_i\right) \setminus Z = 
 \emptyset.
\end{eqnarray}
For each $i$, $Z_i$ is a patch closed subset of $S$, as is $Z$. For each projective model $Y$ of $F/D$, the map $d_Y$ is continuous  in the patch topology {\bf [ref]}, and since $Y$ is Hausdorff in the patch topology, this map is closed in the patch topology. Thus $Y(Z_i) = d_Y(Z_i)$ is patch closed, as is $Y(Z)$.  Since $Y$ is a Noetherian space, the Zariski closed sets in $Y$ are   patch clopen sets in $Y$. Since $d_Y$ is closed in the Zariski topology,  $Y(Z)$ is a patch clopen set in $Y$. Thus for each $i$, $Y(Z_{i}) \setminus Y(Z)$ is a patch closed subset of $Y$. Consequently, 
$d_Y^{-1}\left( Y(Z_i) \setminus Y(Z)\right)$ is a patch closed set in $\Zar_F(D)$. 

These considerations, along with the fact that $\Zar(D)$ is quasicompact in the patch topology {\bf [ref]}, imply 

there are indices $i_1,\ldots,i_n$  such that
$$\bigcap_{k}d^{-1}\left( d(Z_{i_k}) \setminus d(Z)\right) = \emptyset.$$


Let $Y$ be the join of $Y_1,\ldots,Y_m$. Then 
\begin{center} $d_{Y}^{-1}\left(Y(Z_{i_k})\setminus Y(Z)\right) \: \subseteq \: d_{Y_j}^{-1}(Y_j(Z_{i_k})\setminus Y_j(Z))$ \:for all $j,k$,
\end{center}
and so 
\begin{center} $d_{Y}^{-1}(\bigcap_{k}Y(Z_{i_k})\setminus Y(Z)) = \emptyset$
\end{center}
Since $d_Y$ is surjective, this implies $\bigcap_{k}Y(Z_{i_k})= Y(Z).$ Since $\{D_i\}$ is a directed set, we may choose $i$ such that $D_{i_1},\ldots,D_{i_n} \subseteq D_i.$ For each $k$, $D_i$ dominates $D_{i_k}$ and so  $$Z_i = \Dom_F(D_i) \subseteq \Dom_F(D_{i_k}) = Z_{i_k}.$$ 
Thus $$d_Y^{-1}(Y(Z_i) \setminus Y(Z)) \: \subseteq \: \bigcap_{k}
d_Y^{-1}(Y(Z_{i_k}) \setminus Y(Z_{i})) = \emptyset.$$ It follows that $Y(Z_i) = Y(Z)$, and since $Y$ dominates $X$, we conclude $X(Z_i) = X(Z).$ Since $d_X$ is a closed map, the Zariski closure of $X(Z_i)$ is $X(\overline{Z_i})$ and the Zariski closure of $X(Z)$ is $X(\overline{Z})$.
We have established that  $Z_i$ is connected, and so  since continuous images of connected sets are connected,  $X(Z_i) = X(Z)$ is connected. 
Thus $X(\overline{Z})$ is connected, which proves  $\overline{Z}$ is connected by Lemma~\ref{connected iff}.
\end{proof}




If ${Z}$ is connected, then $X({Z})$ is connected for each projective model $X$ of $F/D$  since a continuous image of a connected set is connected. Conversely, suppose  $X({Z})$ is connected for each projective model $X$ of $F/D$. To show  ${Z}$ is connected, 
suppose by way of contradiction that ${Z}$ is a disjoint union of two nonempty closed sets $ {Z}_1 $ and ${Z}_2$.  For each projective model $X$ of $F/D$, 
$X(Z_1)$ and $X(Z_2)$ are closed sets in $X$ since the domination mapping $d_X:\Zar_F(D) \rightarrow X$ is a closed map. 
Thus since $X({Z})$ is connected  and
$X({Z}) = X({Z}_1) \cup X({Z}_2)$,  it follows that $X({Z}_1) \cap X({Z}_2)$ is nonempty. Let $d_X:\Zar_F(D) \rightarrow X$ denote the domination map that sends a valuation ring in $\Zar_F(D)$ to its center. 
For each projective model $X$, define a closed subset $Z_X$ of $Z$ by    $${Z}_X:= Z \cap  d_X^{-1}(X({Z}_1) \cap X({Z}_2)).$$
The set $Z_X$ is nonempty since  $X({Z}_1) \cap X({Z}_2)$ is nonempty and $Z =Z_1 \cup Z_2$.  
%Then $  {Z}_X \subseteq  {Z}$ for all choices of $X$. 

We claim that  if $X_1,\ldots,X_n$ are projective models of $F/D$, then the intersection of the ${Z}_{X_i}$ is nonempty. Let $X$ be the join of the $X_i$, as in
\cite[p.~120]{MR0389876}. 
  Then $X$ dominates each $X_i$, and the image of $X({Z}_1)$ in $X_i$ under the domination map $d_{X,X_i}:X \rightarrow X_i$ is $X_i({Z}_1)$. Similarly, the image of $X({Z}_2)$ is $X_i({Z}_2)$. 
If $V \in {Z}_X$, then $V$ dominates a local ring in 
$X({Z}_1) \cap X({Z}_2)$, and hence also in $X_i({Z}_1) \cap X_i({Z}_2)$ for each $i$. 
Thus $V \in Z_{X_i}$ and it 
 follows that ${Z}_X \subseteq {Z}_{X_1} \cap \cdots \cap {Z}_{X_n}.$
Since  ${Z}_{X}$ is nonempty, we conclude ${Z}_{X_1} \cap \cdots \cap  {Z}_{X_n}$ is nonempty. This proves the collection   $\{{Z}_X:X$ is a projective model$\}$ of  closed subsets ${Z}_X$  has the finite intersection property. 
Since ${Z}$ is quasicompact (it is a closed subset of the quasicompact space Zar($F/D$)), there is a valuation ring $V \in \bigcap_{X}{Z}_X$.
Since ${Z}$ is a disjoint union of ${Z}_1$ and $ {Z}_2$, we can 
assume without loss of generality that
 $V \not \in  {Z}_2$. 
 
 For each projective model $X$ of $F/D$, since $d_X(V) \in X({Z}_1) \cap X({Z}_2)$, there is a valuation ring $V_X \in {Z}_2$ such that $d_X(V_X) = d_X(V)$. We claim $V$ is a Zariski limit point of the $V_X$. Let ${\mathcal{U}}$ be an open neighborhood of $V$.  There are $x_1,\ldots,x_m \in F$ such that $V \in  {\mathcal{U}}(x_1,\ldots,x_m) \subseteq {\mathcal{U}}.$ The projective model $X$ of $F/D$ defined by $1,x_1,\ldots,x_m$ has the property that the the localizations of $D[x_1,\ldots,x_n]$ at prime ideals are all in $X$.  Since $d_X(V)$ is a localization of $D[x_1,\ldots,x_n]$ at a prime ideal and  $d_X(V) = d_X(V_X)$, we have $D[x_1,\ldots,x_n] \subseteq V_X$, so that $V_X \in {\mathcal{U}}(x_1,\ldots,x_m)$. 
 This shows
 $V$ is a limit point of the $V_X$, where $X$ ranges over the projective models of $F/D$.   But
 the $V_X$ are all in the closed set ${Z}_2$ while the limit point $V$ is not. This contradiction shows ${Z}$ is connected. 
 
 
 
\begin{lemma} \label{closure} Let $D$ be a subring of a field $F$, and let $Z$ be a nonempty patch closed subset of $\Zar_F(D)$. Then $Z= \bigcap_Yd_Y^{-1}(Y(Z))$, where $Y$ ranges over all the projective models  of $F/D$.  
\end{lemma}

\begin{proof}
%It is clear that $Z \subseteq \bigcap_Yd_Y^{-1}(Y(Z))$. 
Let $V \in \bigcap_Yd_Y^{-1}(Y(Z))$. 
We show $V$ is a patch limit point of $Z$. Let $x_1,\ldots,x_n \in V$, and let $y \in F \setminus V$.  Let $Y$ be the projective model defined by $x_1,\ldots,x_n,y^{-1}$. Since $V \in d^{-1}_Y(Y(Z))$, there is a valuation ring $V_Y \in Z$ such that $d_Y(V) = d_Y(V_Y)$. Since $x_1,\ldots,x_n,y^{-1} \subseteq V$, the rings $V$ and $V_Y$ dominate a localization of $D[x_1,\ldots,x_n,y^{-1}]$. 
Since $y^{-1} \in {\mathfrak{M}}_V$,  it follows that $y^{-1}  \in  {\mathfrak{M}}_{V_Y}$. Thus $V$ is a patch limit point of the $V_Y$, and so $V$ is a patch limit point of $Z$. Since $Z$ is patch closed, this shows $V \in Z$.  
\end{proof}


****************


\begin{theorem} \label{Dom(R)} Let $R$ be an integrally closed local domain with quotient field $F$. 
 Suppose there is a Noetherian subring $D$ of $R$ whose quotient field is separably closed in $F$, and suppose 
$R$ is a directed limit of normal Noetherian   $D$-subalgebras $D_i$ of $R$. Then the Zariski closure of $\Dom_F(R)$ in $\Dom_F(D)$ is connected.   
\end{theorem} 

\begin{proof}  
By assumption, we may 
 write $R = \bigcup_i D_i$ for a directed set of normal Noetherian  local $D$-subalgebras of $R$, each of which is dominated by $R$. 
%Since $D$ is analytically unramified we may write $A(\overline{Z}) = \bigcup_i D_i$ for a directed set of normal Noetherian local $D$-subalgebras of $A(\overline{Z})$, each of which is dominated by $A(\overline{Z})$ and has quotient field $F$. 
Let $Z$ denote the Zariski closure of $\Dom_F(R)$ in $\Dom_F(D)$, and for each $i$, let $Z_i$ denote the Zariski closure of $\Dom_F(D_i)$ in $\Dom_F(D)$.


  We first claim each $ {Z_i}$ is connected. By Lemma~\ref{connected iff}, it suffices to show $X(Z_i)$ is connected for each projective model $X$ of $F/D$. Fix an $i$ in the index set of the $Z_i$, and let $X$ be a projective model of $F/D$. 
By {\bf [ref]}, the image of $\Dom_F(D_i)$ in $X$ under the domination map $d_X$ is the fiber of
 the map $X \rightarrow \Spec D_i$ over the closed point 
  of $\Spec D_i$. Since the quotient field of $D$ is separably closed in $F$, 
the quotient field of $D$ is also separably closed in the   
   function field of $X$. 
  Thus, by Zariski's Connectedness Theorem, the fact that $D_i$ is normal, Noetherian and local implies the fiber of the map
  $X \rightarrow \Spec D_i$ over the closed point 
  of $\Spec(D_i)$ is
     connected in the Zariski topology.
  Since $d_X$ is a closed map,   the image $X( {Z_i})$ of $Z_i$ under $d_X$ is the Zariski closure of this fiber $X(\Dom_F(D_i))$ in $X$.  So $X({Z_i})$ is connected since $X(\Dom_F(D_i))$ is connected and dense in $X( {Z_i})$. 
  Thus
    ${Z_i}$ is connected.

Now to show $Z$ is connected it is enough by Lemma~\ref{spectral connected 2} to show $Z = \bigcap_i Z_i$.  (In applying the lemma, $S$ is $\Zar_F(D)$ and the $X_\alpha$ are the projective models of $F/D$.) 
  Let $V \in \bigcap_i  {Z}_i$. By Lemma {\bf ref}, we have for each $i$ that  there is $U_i \in \Dom_F(D_i)$ such that $V \subseteq U_i$. As overrings of $V$, the rings in $\{U_i\}$ are comparable with respect to set inclusion, and so $U:=\bigcup_i U_i$ is a valuation ring. Now $R = \bigcup_i D_i \subseteq \bigcup_i U_i = U$, and ${\mathfrak m}_R = \bigcup_{i} {\mathfrak{m}}_{D_i} \subseteq \bigcup_i {\mathfrak{m}}_{U_i} = {\mathfrak m}_U$. Thus $U \in \Dom_F(R)$, and since $D \subseteq V \subseteq U$, we have $V \in  {Z}$ by {\bf ref}. This verifies that $ {Z} = \bigcap_{i}  {Z}_i$. %and this completes the proof.
  
  Finally, we show $\Dom_F(R)$ is connected. Suppose $\Dom_F(R) = Z_1 \cup Z_2$ for disjoint closed subsets $Z_1$ and $Z_2$ of $\Dom_F(R)$.  Let $\overline{Z_1}$ and $\overline{Z_2}$ be the Zariski closures of $Z_1$ and $Z_2$ in $\Dom_F(D)$. Then $Z = \overline{Z_1} \cup \overline{Z_2}$.
  Since $Z_1$ and $Z_2$ are Zariski closed in $\Dom_F(R)$ and $\Dom_F(R)$ is patch closed in $\Dom_F(D)$, the sets $Z_1$ and $Z_2$ are patch closed in $\Dom_F(D)$, and so $\overline{Z_i}$ is the set of valuation rings in $\Dom_F(D)$ that are  contained in valuation rings in $Z_i$. 
  Since $Z$ is connected, there is $V \in \overline{Z_1} \cap \overline{Z_2}$.  Thus there are $U_1 \in Z_1$ and $U_2 \in Z_2$ such that $V \subseteq U_1 \cap U_2$.
 The overrings of a valuation ring are totally ordered under inclusion, so  
  without loss of generality, $U_1 \subseteq U_2$.  Now $U_1,U_2 \in \Dom_F(R)$ and $Z_2$ is closed in  $\Dom_F(R)$, so  $U_1 \in Z_2$ since $U_1 \subseteq U_2 \in Z_2$. However, by assumption, $Z_1 \cap Z_2 = \emptyset$. This contradiction shows $\Dom_F(R)$ is connected.    
\end{proof} 


**************


\begin{theorem} \label{Dom(R)} Let  $R$ be a local subring of a field $F$ such that  
$R$ is a directed union of 
 Noetherian  local subrings $D_i$ of $R$, each of which is dominated by $R$ and integrally closed in $F$.
Then $\Dom_F(R)$ is Zariski connected.    
\end{theorem} 

\begin{proof} 
%Let $D$ be the localization of the prime subring of $R$ at the prime ideal that is contracted from the maximal ideal of $R$. 
%Let $\{S_\alpha\}$ be the collection of finite subsets of $R$. Since $D$ is an excellent local domain, the integral closure $B_\alpha$ of $D$ in the quotient field $F_\alpha$ of $D[S_\alpha]$ is module-finite and hence Noetherian. Since $R$ is integrally closed in $F$ and $D \subseteq R$, we have that $B_\alpha \subseteq R$. Let $D_\alpha$ be the localization of $B_\alpha$ at the prime ideal that is contracted from  the maximal ideal of $R$ to $B_\alpha$. Then $R$ is a directed union of the $D_\alpha$.  
%
%For each $\alpha$, let $X_\alpha$ be the projective model of $F/D_\alpha$ defined by $S_\alpha$. Then $\Zar_F(R)$ is the inverse limit of the $X_\alpha$ since each projective model of $F/D$ is dominated by a projective model of $F/D_\alpha$, which in turn is dominated by $\Zar_F(R)$.  

%Let $Z = \Zar_F(R)$, and for each $\alpha$, let $Z_


%Let $\{B_i\}$ be the collection of finitely generated $B$-subalgebras of $R$, and for each $i$ denote by $F_i$ the quotient field of $B_i$. Since $B$ is an excellent Noetherian local domain and $B_i$ is a finitely generated $B$-algebra, the integral closure of $B_i$ in $F_i$ is Noetherian and is a subring of $R$. Let $D_i$ be the localization of this ring at the prime  ideal  that is contracted from the maximal ideal of $R$. Then $R$ is a directed union of the Noetherian local domains $D_i$, with each $D_i$   integrally closed in $F_i$.  





%The integral closure  in $F$ of  a finitely generated $D$-subalgebra of $R$ is module-finite since $D$ is excellent, and so each such integral closure is Noetherian and integrally closed in $F$. This remains true for the localization of the subalgebra at the prime ideal that is contracted from the maximal ideal of $R$. In this way we may represent
%Write $R$ as a directed union $R=\bigcup_i D_i$ of 
% Noetherian  local subrings $D_i$ of $R$, each of which is dominated by $R$ and integrally closed in $F$.  
%Since $D$ is analytically unramified we may write $A(\overline{Z}) = \bigcup_i D_i$ for a directed set of normal Noetherian local $D$-subalgebras of $A(\overline{Z})$, each of which is dominated by $A(\overline{Z})$ and has quotient field $F$. 
%Since $R$ is a directed union of the $D_i$, we may assume
%without loss of generality that among the rings $D_i$ there is one, $D$, that is contained in all the other $D_i$. 
%
%Let $A$ be the localization of the integral closure of $D$ in $F$ at the prime ideal contracted from the maximal ideal of $R$.  
%Show $\Dom_F(A)$ is connected. 
Let $Z$ denote the Zariski closure of $\Dom_F(R)$ in $\Dom_F(D)$, 
and for each $\alpha$, let $Z_\alpha$ denote the Zariski closure of $\Dom_F(D_\alpha)$ in $\Dom_F(D)$.
  We first claim each $ {Z_\alpha}$ is connected. By Lemma~\ref{connected iff}, it suffices to show $X(Z_\alpha)$ is connected for each projective model $X$ of $F/D_i$. Fix an $i$ in the index set of the $Z_i$, and let $X$ be a projective model of $F/D_i$. 
By {\bf [ref]}, the image of $\Dom_F(D_i)$ in $X$ under the domination map $d_X$ is the fiber of
 the map $X \rightarrow \Spec D_i$ over the closed point 
  of $\Spec D_i$. Since  $D_i$ is integrally  closed in $F$, 
$D_i$ is integrally closed in the   
   function field of $X$. 
  Thus, by the version of Zariski's Connectedness Theorem discussed in the introduction, the fact that $D_i$ is Noetherian, local and integrally closed in $F$  implies the fiber of the map
  $X \rightarrow \Spec D_i$ over the closed point 
  of $\Spec D_i$ is
     connected in the Zariski topology.
  Since $d_X$ is a closed map,   the image $X( {Z_i})$ of $Z_i$ under $d_X$ is the Zariski closure of  $X(\Dom_F(D_i))$ in $X$.  So $X({Z_i})$ is connected since $X(\Dom_F(D_i))$ is connected and dense in $X( {Z_i})$. 
  Thus
    ${Z_i}$ is connected.

Now to show $Z$ is connected it is enough by Lemma~\ref{spectral connected 2} to show $Z = \bigcap_i Z_i$.  (In applying the lemma, $S$ is $\Zar_F(D)$ and the $X_\alpha$ are the projective models of $F/D$.) 
It is clear that $Z \subseteq \bigcap_i Z_i$ since $\Dom_F(R) \subseteq \Dom_F(D_i)$ for each $i$. To verify the reverse inclusion,   let $V \in \bigcap_i  {Z}_i$. By Lemma~\ref{generic closure}, we have for each $i$ that  there is $U_i \in \Dom_F(D_i)$ such that $V \subseteq U_i$. As overrings of $V$, the rings in $\{U_i\}$ are comparable with respect to set inclusion, and so $U:=\bigcup_i U_i$ is a valuation ring. Now $R = \bigcup_i D_i \subseteq \bigcup_i U_i = U$, and ${\mathfrak m}_R = \bigcup_{i} {\mathfrak{m}}_{D_i} \subseteq \bigcup_i {\mathfrak{m}}_{U_i} = {\mathfrak m}_U$. Thus $U \in \Dom_F(R)$, and since $D \subseteq V \subseteq U$, we have $V \in  {Z}$ by Lemma~\ref{generic closure}. This verifies that $ {Z} = \bigcap_{i}  {Z}_i$, and hence $Z$ is connected. To complete the proof,  observe that since $\Dom_F(R)$ is patch closed in $\Dom_F(D)$, hence also in $\Zar_F(D)$, and $\Dom_F(R)$ is Zariski dense in $Z$, Lemma~\ref{dense} shows $\Dom_F(R)$ is connected. 
%
%and this completes the proof.
 % 
  %Finally, we show $\Dom_F(R)$ is connected. Fix $i$ and write $D=D_i$. Suppose $\Dom_F(R) = Z_1 \cup Z_2$ for disjoint closed subsets $Z_1$ and $Z_2$ of $\Dom_F(R)$.  Let $\overline{Z_1}$ and $\overline{Z_2}$ be the Zariski closures of $Z_1$ and $Z_2$ in $\Dom_F(D)$. Then $Z = \overline{Z_1} \cup \overline{Z_2}$.
 % Since $Z_1$ and $Z_2$ are Zariski closed in $\Dom_F(R)$ and $\Dom_F(R)$ is patch closed in $\Dom_F(D)$, the sets $Z_1$ and $Z_2$ are patch closed in $\Dom_F(D)$, and so $\overline{Z_i}$ is the set of valuation rings in $\Dom_F(D)$ that are  contained in valuation rings in $Z_i$. 
 % Since $Z$ is connected, there is $V \in \overline{Z_1} \cap \overline{Z_2}$.  Thus there are $U_1 \in Z_1$ and $U_2 \in Z_2$ such that $V \subseteq U_1 \cap U_2$.
% The overrings of a valuation ring are totally ordered under inclusion, so  
 % without loss of generality, $U_1 \subseteq U_2$.  Now $U_1,U_2 \in \Dom_F(R)$ and $Z_2$ is closed in  $\Dom_F(R)$, so  $U_1 \in Z_2$ since $U_1 \subseteq U_2 \in Z_2$. However, by assumption, $Z_1 \cap Z_2 = \emptyset$. This contradiction shows $\Dom_F(R)$ is connected.    
\end{proof} 


***************

\begin{corollary} \label{Nagata ring} 
Let $D$ be analytically unramified Noetherian local domain. If $R$ is a local integrally closed overring of $D$, then $\Dom_F(R)$ is connected. 
%
%Let $R$ be an integrally closed local domain having quotient field $F$. If there is a analytically unramified Noetherian local subring $D$ of $R$ such  that $D$ has the same residue field as $R$, $D$   is a Nagata ring and the  quotient field of $D$ is separably closed in $F$, then $\Dom_F(R)$ is connected. 



\end{corollary}




\begin{proof} 
%Let $Z$ denote the Zariski closure of $\Dom_F(R)$ in $\Zar_F(D)$.
%We claim first that $R = A(Z)$. 
Since $D$ is analytically unramified,   the integral closure of each finitely generated $D$-subalgebra of $R$ in $F$ is again Noetherian. Thus Theorem~\ref{Dom(R)} implies $\Dom_F(R)$ is connected. 
\end{proof} 

\begin{corollary} \label{au} Let $D$ be analytically unramified Noetherian local domain with quotient field $F$. If $R$ is a local integrally closed overring of $D$, then for each   projective model $X$ of $F/R$, the fiber of the map $d:X \rightarrow \Spec R$ over the closed point of $\Spec R$ is connected. 
\end{corollary} 

\begin{proof}
The fiber of $d$ over the closed point of $\Spec R$ is a continuous image of $\Dom_F(R)$, so the corollary follows from Theorem~\ref{Nagata ring}.
\end{proof}

**************

A map $d:X \rightarrow Y$ of spectral spaces is {\it spectral} if it is continuous in the Zariski and patch topologies.  The category of spectral spaces has as its objects spectral spaces and as its morphisms spectral maps. 
 
 {\bf THE FOLLOWING LEMMA HAS A GAP IN THE PROOF AND SO MAY NOT BE CORRECT}
 
 \begin{lemma} \label{spectral connected 2} 
 Let $S$ be a spectral space that is the inverse limit of Noetherian spectral spaces $X_\alpha$ in the category of spectral spaces, and suppose   each projection map $d_\alpha:S \rightarrow X_\alpha$ is   closed. 
 If  $\{Z_i\}$ is an inverse system  of closed connected subsets of $S$, then
 $\bigcap_{i}Z_i$ is connected. 
 \end{lemma}
 
 \begin{proof} 
 Let  $Z = \bigcap_{i}Z_i$.
Then $Z$ is closed and hence quasicompact, so to prove $Z$ is connected, it suffices by Lemma~\ref{spectral connected 1} to show $d_\gamma(Z)$ is connected for each $\gamma$ in the index set of the $Z_\alpha$. Fix $\gamma$. 
 Since $Z$ and the $Z_i$ are   closed,  Lemma~\ref{easy} %and the fact that $S$    is also an inverse limit of the $X_\alpha$ with respect to the patch topology on $S$ and $X_\alpha$ {\bf [ref]}, 
 implies
\begin{center} 
$Z = \bigcap_{\alpha \geq \gamma} d_\alpha^{-1}(d_\alpha(Z))$ \: and \: $Z_i =  \bigcap_{\alpha \geq \gamma} d_\alpha^{-1}(d_\alpha(Z_i))$ for all $i$.\end{center}  
%Since the image of a connected set under a continuous map is connected, to show $X(Z)$ is connected it suffices to show 
% $X(Z) = X(Z_i)$ for some $i$.
Thus
%For each $i$ and $\alpha$, $d_\alpha(Z) \subseteq d_\alpha(Z_i)$ and $$d_\alpha^{-1}(d_\alpha(Z_i) \setminus d_\alpha(Z)) = d_\alpha^{-1}(d_\alpha(Z_i)) \setminus d_\alpha^{-1}(d_\alpha(Z))$$
% and we have 
\begin{eqnarray} \label{eq d}
\bigcap_i\bigcap_{\alpha \geq \gamma}d_\alpha^{-1}\left( d_\alpha(Z_i) \setminus d_\alpha({Z})\right)
& \subseteq  & \bigcap_i \left(
\bigcap_{{\alpha \geq \gamma}}d_\alpha^{-1}(d_\alpha(Z_i))\right)\\
&= &  \bigcap_i Z_i \: \: = \:\: Z.
\end{eqnarray}
However, it is clear that $$Z \cap \left(\bigcap_i\bigcap_{\alpha \geq \gamma}d_\alpha^{-1}\left( d_\alpha(Z_i) \setminus d_\alpha({Z})\right)\right) = \emptyset,$$
and so 
\begin{eqnarray}\label{empty eq} 
\bigcap_i\bigcap_{\alpha \geq \gamma}d_\alpha^{-1}\left( d_\alpha(Z_i) \setminus d_\alpha({Z})\right) & = & \emptyset.
\end{eqnarray}



%\bigcap_i \bigcap_{d_\alpha}d_Y^{-1}( %Y(Z_i) \setminus Y(Z)) 
%%& \subseteq & \bigcap_i \left(
%%\bigcap_{{\alpha \geq \gamma}}d_\alpha^{-1}(d_\alpha(Z_i))\right)\setminus \bigcap_{{\alpha \geq \gamma}}d_\alpha^{-1}(d_\alpha({Z})) \\ 
%%&= & \left(\bigcap_i Z_i\right) \setminus {Z} = %%
%% \emptyset.
%%\end{eqnarray}
 For each ${\alpha \geq \gamma}$, since the map $d_\alpha$ is closed, $d_\alpha(Z)$ and the  $d_\alpha(Z_i)$ are closed in $X_\alpha$.     Since $X_\alpha$ is a Noetherian space, the
 open subsets of $X_\alpha$ are quasicompact and hence patch clopen, and thus the closed subsets of $X_\alpha$ are patch clopen also. Consequently, 
 %$d_\alpha({Z})$ is a patch clopen set in $X_\alpha$, and so 
 for each $i$, $d_\alpha(Z_{i}) \setminus d_\alpha( {Z})$ is a patch closed subset of $X_\alpha$. Thus since the maps $d_\alpha$ are spectral, 
$d_\alpha^{-1}\left( d_\alpha(Z_i) \setminus d_\alpha({Z})\right)$ is a patch closed set in $S$    for each ${\alpha \geq \gamma}$.


 
Since  $S$ is quasicompact in the patch topology and the sets in the intersection in equation~(\ref{empty eq}) are patch closed, we conclude   there are $\alpha_1,\ldots,\alpha_m \geq \gamma$ and indices $i_1,\ldots,i_n$  such that
$$\bigcap_{j,k}d_{\alpha_j}^{-1}\left( d_{\alpha_j}(Z_{i_k}) \setminus d_{\alpha_j}({Z})\right) = \emptyset.$$ 
Let $\beta \geq \alpha_1,\ldots,\alpha_m$. 
Now {\bf THE FOLLOWING CONTAINMENT MAY BE INCORRECT. IF IT'S CORRECT IT NEEDS A PROOF}




\begin{center} $d_{\beta}^{-1}\left(d_\beta(Z_{i_k})\setminus d_\beta({Z})\right) \: \subseteq \: d_{\alpha_j}^{-1}(d_{\alpha_j}(Z_{i_k})\setminus d_{\alpha_j}({Z}))$ \:for all $j,k$,
\end{center}
and so 
\begin{center} $d_{\beta}^{-1}(\bigcap_{k}d_\beta(Z_{i_k})\setminus d_\beta({Z}))=\emptyset$
\end{center}
Since $d_\beta$ is surjective, this implies $\bigcap_{k}d_\beta(Z_{i_k}) \subseteq d_\beta({Z}).$ 
Now since $\{Z_i\}$ is an inverse system, we may choose $i$ such that $Z_i \subseteq Z_{i_1} \cap \ldots \cap Z_{i_n},$ and so 
 $d_\beta(Z_{i}) \subseteq d_\beta({Z})$. Since $Z \subseteq 
 Z_i$, we conclude $d_\beta(Z_{i})=d_\beta(Z)$. By assumption, $Z_i$ is connected, so $d_\beta(Z_i) = d_\beta(Z)$ is connected, and since $\beta \geq \gamma$, $d_\gamma(Z)$, as a continuous image of $d_\beta(Z)$, is connected. The choice of $\alpha$ was arbitrary, so  $Z$ is connected. 
 \end{proof}


 
\begin{lemma}
If $R$ is a local ring that is integrally closed in a field $F$, and $D$ is a local subring of $R$ that is dominated by $R$ and has the same residue field as $R$, then $R = A(Z)$, where $Z$ is the Zariski closure of $\Dom_F(R)$ in $\Dom_F(D)$. 
\end{lemma} 

\begin{proof} 
Since $R$ is local and has the same residue field as $D$, we have $R = D+{\mathfrak{m}}_R$. Since $R$ is integrally closed, ${\mathfrak{m}}_R = J(\Dom_F(R))$.   
 By Lemma~\ref{J lemma}, $J(
 \Dom_F(R)) = J({Z})$, and so $R =  D+{\mathfrak{m}}_R = D+J({Z}) \subseteq A({Z}) \subseteq A(\Dom_F(R)) = R$, proving that $R = A({Z})$.
 \end{proof} 

 We have $D+{\mathfrak{m}}_R \subseteq $A(Z)$ since ${\mathfrak{m}}_R = J(\Dom_F(R)) = J(Z)$. 

 Let $U \in Z$ and $V \in \Dom_F(D)$ such that $V \subseteq U$.  Then ${\mathfrak{m}}_R \subseteq {\mathfrak{m}}_V$.  Suppose $R \not \subseteq V$.  Then $x \in R$ with $x^{-1} \in {\mathfrak{m}}_V$. 



 So if $V$ contains $D+{\mathfrak{m}}_R$, then $V$ contains $R. If $V$ 

 Let $U \in Z$ and $V \in \Dom_F(D)$ such that $V \subseteq U$.  Then ${\mathfrak{m}}_R \subseteq {\mathfrak{m}}_V$ so $R \subseteq V$.

 ***********

 \begin{lemma} \label{Noetherian local} If $D$ is     local subring of a field $F$ such that $D$ is integrally closed in $F$,  
 then $\Dom_F(D)$ is    connected.  
\end{lemma} 



\begin{proof} Let $Z=\Dom_F(D)$. To show $Z$ is connected, it suffices by  Lemma~\ref{connected iff} to show $X(Z )$ is connected for each projective model $X$ of $F/D$. Let $X$ be a projective model of $F/D_i$. 
As follows from the property of completeness discussed in Section~2 (or from the valuative criterion for properness), the image $X(Z)$ of $Z$ in $X$ under the domination map $d_X$ is the fiber of
 the map $X \rightarrow \Spec D$ over the closed point 
  of $\Spec D$. Since  $D$ is integrally  closed in $F$, 
$D$ is integrally closed in the   
   function field of $X$. 
  Thus, by the version of Stein factorization   discussed in the introduction, the fact that $D$ is  local and integrally closed in $F$  implies the fiber $X(Z)$ of the map
  $X \rightarrow \Spec D$ over the closed point 
  of $\Spec D$ is
     connected in the Zariski topology.
%  Since $d_X$ is a closed map,   the image $X( {Z})$ of $Z$ under $d_X$ is the Zariski closure of  $X(Z)$ in $X$.  
This proves
    ${Z}$ is connected.
 \end{proof} 


 DIDN"T NEED:


 {\bf THE NEXT  LEMMA IS NOT USED ANYWHERE}

\begin{lemma} \label{up} Let $D$ be a local subring of a field $F$,  let $Z$ be a patch closed subset of $\Dom_F(D)$ and let $\overline{Z}$ be the Zariski closure  of $Z$ in $\Dom_F(D)$.
%\begin{enumerate}
%\item[$(1)$] If $\overline{Z}$  is inversely connected, then $Z$ is inversely connected. 
 If  $A(Z)$ is a  local ring, then $A(\overline{Z})$ is a local ring.
\end{lemma} 

\begin{proof} 
%Suppose $\overline{Z}$ is connected, and let $Z_1$ and $Z_2$ be Zariski closed subsets of $Z$ such that $Z = Z_1 \cup Z_2$. Then $\overline{Z} = \overline{Z_1} \cup \overline{Z_2}$. Since $\overline{Z}$ is connected, there is a valuation ring $V \in \overline{Z_1} \cap \overline{Z_2}$.  
%Since $Z$ is patch closed in $\Zar_F(D)$, so are $Z_1$ and $Z_2$. Thus by {\bf [ref]}, there are valuations ring $U_1 \in Z_1$ and $U_2 \in Z_2$ such that $V \subseteq U_1 \cap U_2$. The valuation overrings of a valuation ring are comparable under set inclusion, so 
 Let $x$ and $y$ be nonunits in $A(\overline{Z})$, and suppose by way of contradiction that $x+y$ is a unit in $A(\overline{Z})$. Then $x+y$ is a unit in $A(Z)$ and since $A(Z)$ is local, then implies at least one of $x$ and $y$ is a unit in $A({Z})$, say $x$. Since $x$ is a nonunit in $A(\overline{Z})$, there is a valuation ring $V \in \overline{Z}$ such that $x \in {\mathfrak{M}}_V$. By {\bf [ref]}, $V \subseteq U$ for some valuation ring $U \in Z$.  All of the valuation rings contained in $U$ and dominating $D$ are in $\overline{Z}$. Since $D$ is local, {\bf [ref]} implies the intersection of all these rings is local with maximal ideal ${\mathfrak M}_U$. Since $x$ is a nonunit in this intersection, $x$ is
in ${\mathfrak M}_U$. But then since $A(Z) \subseteq U$,  $x$ is a nonunit in $A(Z)$, a contradiction that implies $A(\overline{Z})$ is local. 
\end{proof}

